\documentclass[11pt]{article}
\usepackage[margin=1in]{geometry}
\usepackage{amsmath,amssymb,amsthm}
\usepackage{fontspec}
\setmainfont{Latin Modern Roman}
\setsansfont{Latin Modern Sans}
\setmonofont{DejaVu Sans Mono}[Scale=0.86]
\usepackage{xcolor}
\usepackage{booktabs}
\usepackage{array}
\usepackage{calc}
\usepackage{longtable}
\usepackage{seqsplit}
\usepackage{fancyvrb}
\usepackage{fvextra}
\fvset{breaklines=true,breakanywhere=true,fontsize=\small}
\usepackage[hidelinks]{hyperref}
\usepackage{url}

\providecommand{\textquotesingle}{'}
\providecommand{\textquotedbl}{"}
\providecommand{\tightlist}{\setlength{\itemsep}{0pt}\setlength{\parskip}{0pt}}

\newtheorem{theorem}{Theorem}
\newtheorem{lemma}[theorem]{Lemma}
\newtheorem{proposition}[theorem]{Proposition}
\theoremstyle{definition}
\newtheorem{definition}[theorem]{Definition}
\theoremstyle{remark}
\newtheorem{remark}[theorem]{Remark}

% A declaration of the development; the name/line gate (layer 11) resolves every File:line anchor
% against the pinned tree.
\DeclareUrlCommand\decl{\urlstyle{tt}\def\UrlBreaks{\do\_\do\.}}
\newcommand{\plcommit}{d15464f0f4f4c8848e7d3e63dd68a03ce45433ae}
\newcommand{\src}[2]{\href{https://github.com/jpbald93/paperl_clifford_chain/blob/d15464f0f4f4c8848e7d3e63dd68a03ce45433ae/#1\#L#2}{\texttt{\seqsplit{#1}}:#2}}
\newcommand{\PH}[1]{\textcolor{red}{\textbf{[PLACEHOLDER:} #1\textbf{]}}}

\title{Clifford's chain for every $n$: point--circle configurations on the Riemann sphere, in Lean~4}
\author{Josh Bald\\ \small Independent researcher, Ontario, Canada\\ \small\texttt{jpbald93@gmail.com}\\
\small ORCID: \texttt{0009-0002-1317-6489}
\thanks{\copyright~2026 Josh Bald. Licensed under CC BY 4.0
(\url{https://creativecommons.org/licenses/by/4.0/}). The accompanying Lean code is licensed under
the Apache License 2.0.}}
\date{26 November 2026}
\emergencystretch=3em
\begin{document}
\sloppy
\maketitle

\begin{abstract}
We report a Lean 4 formalisation, on top of Mathlib, of Clifford's chain of circle
theorems for every \(n\), on the Riemann sphere rather than in the finite plane. For every family of
\(n\) circles through a common point we prove that a level-\(n\) Clifford object exists: for odd \(n\) a
generalized circle (lines included) meeting one level-\((n-1)\) object of each deletion, for even \(n\) a
point of \(\mathrm{OnePoint}\,\mathbb{C}\) ($\infty$ included) lying on one such object of each deletion. The
odd side is proved without hypotheses beyond general position; the even side, including the
degenerate corner where the parent cross vector vanishes, is proved in full. Base levels are
transported from the finite-plane theorem, and a six-centre configuration shows that the finite plane
itself is false for \(n \ge 6\), which is why the sphere is needed. All terminal theorems depend only
on Lean's standard axioms. The paper reports the verification record in full --- an axiom gate over 36
declarations, an eight-test tampering suite, mutation testing over 516 of the 577 theorem-level
declarations, contradiction probes over 551, a 262-declaration hypothesis ablation, and
satisfiability certification on two machines --- and states its limits: one open lemma in the direct S1
route, four layer-13 rows whose writer-side statements need new mathematics, 180 declarations
unresolved by the satisfiability search, Carver's 1920 criterion unformalised, and the recorded gaps
of the prior-art search.
\end{abstract}

Axiom shorthand used throughout: \textbf{std} = \texttt{{[}propext,\ Classical.choice,\ Quot.sound{]}}.


\hypertarget{introduction}{%
\section{Introduction}\label{introduction}}

Clifford's chain is a family of incidence theorems generated by one recursion. Fix n circles of the
plane through a common point M, in general position. Every three of them have
three second intersections, and those three lie on a generalized circle; every four of them give four
such level-3 circles, and the four share a point; and so on, each level indexed by the subfamilies of
a fixed size. The chain states that the process never dies: at every level the objects either pass
through a common point (even levels) or lie on a common circle (odd levels). The statement and its
first proofs are classical --- Clifford (\emph{Messenger of Mathematics} 5, 1870--71; reprinted in
\emph{Mathematical Papers}, 1882), Morley (1900), Carver (1920), and the polytope picture of
Longuet-Higgins (1972); the Miquel/cross-ratio route is set out by Bobenko and Schief
(arXiv:1509.04109). The appropriate ambient plane is the inversive (M\"obius) plane, in which the
``circle'' of an odd level is allowed to degenerate to a line, and the ``point'' of an even level is
allowed to be $\infty$. Carver's 1920 paper \emph{The Failure of the Clifford Chain} classifies exactly when
this happens: in the inversion plane, with one point at infinity and with circles understood to
include lines and point-circles, every improper set is a limit of proper ones and is detected by a
vanishing minor of a characteristic matrix (his \S{}5, pp.~160--161, with Theorem 13; \S{}6, pp.~161--164,
treats the limiting figures themselves).

\textbf{What this paper contributes is a formalisation, not new mathematics.} The claim under test is
the chain for \emph{all} n on the Riemann sphere, with the level objects alternating between points of
\texttt{OnePoint\ ℂ} ($\infty$ included) and generalized circles (lines included), together with a verification
record for the resulting Lean 4 development. The novelty scout (\texttt{SCOUT.\allowbreak md} \S{}1) records the verdict
\texttt{NONE-FOUND} for the all-n statement in any proof assistant (rows L5 and L6), after searching the
Mathlib tree at pin \texttt{0df444a3} and its \texttt{docs/\allowbreak 1000.\allowbreak yaml}, GitHub code search for Lean and every other
language, the Isabelle AFP, atlas-lean, lean-genius, Rocq/GeoCoq and HighSchoolGeometry, HOL Light,
HOL4, Metamath \texttt{set.\allowbreak mm}, the Mizar MML listing, the formal-conjectures/compfiles/ICML collections,
ProofAtlas, a numeric check of the statement, arXiv, OpenAlex, Crossref and the general web. In
Mathlib, \texttt{Clifford\textquotesingle{}s\ circle\ theorems} appears in \texttt{docs/\allowbreak 1000.\allowbreak yaml} with no \texttt{decl}; the Miquel results
that do exist (AFP \texttt{Miquel.\allowbreak thy}, 2026-08-11, Ramos--Hulak--Queiroz; atlas-lean
\texttt{MathlibExt/\allowbreak Geometry/\allowbreak Euclidean/\allowbreak MiquelTheorem.\allowbreak lean}; Rocq \texttt{HighSchoolGeometry}; lean-genius
\texttt{MiquelPivotTheoremOQ01.\allowbreak lean}) are all in the \emph{pivot} or \emph{complete-quadrilateral} form and are the
base case only. This paper therefore claims \textbf{first machine-checked treatment of the chain for
arbitrary n} --- a claim scoped to the 26 search rows logged in \texttt{SCOUT.\allowbreak md} \S{}0 and to the sources
\S{}5.10 records as not searched --- and nothing about the mathematics.

\hypertarget{what-the-formalisation-closes}{%
\subsection{What the formalisation closes}\label{what-the-formalisation-closes}}

Two statement-level targets were carried through the development:

\begin{itemize}
\tightlist
\item
  \textbf{C5 --- the all-n chain on the sphere.} \texttt{CliffordChainS\ n} for every \texttt{n\ ≥\ 1}. The library states
  this in two pieces: \texttt{cliffordChainS\_\allowbreak of\_\allowbreak le\_\allowbreak five\ (n)\ (hn\ \allowbreak :\ n\ ≤\ 5)} (\href{https://github.com/jpbald93/paperl_clifford_chain/blob/d15464f0f4f4c8848e7d3e63dd68a03ce45433ae/Clifford/Mobius.lean\#L178}{\texttt{Clifford/\allowbreak Mobius.lean}:178}) for
  the low levels, and the headline theorem \texttt{cliffordChainS\_\allowbreak inf\ (n)\ (hn\ \allowbreak :\ 3\ ≤\ n)}
  (\href{https://github.com/jpbald93/paperl_clifford_chain/blob/d15464f0f4f4c8848e7d3e63dd68a03ce45433ae/Clifford/MorleyKernel.lean\#L5592}{\texttt{Clifford/\allowbreak MorleyKernel.lean}:5592}) for \texttt{n\ ≥\ 3}. The chain of calls behind the headline is short:
  \texttt{cliffordChainS\_\allowbreak inf} is \texttt{cliffordChainS\_\allowbreak of\_\allowbreak corner\ EvenCorner\_\allowbreak proof}
  (\href{https://github.com/jpbald93/paperl_clifford_chain/blob/d15464f0f4f4c8848e7d3e63dd68a03ce45433ae/Clifford/MorleyKernel.lean\#L5269}{\texttt{Clifford/\allowbreak MorleyKernel.lean}:5269}), which is \texttt{cliffordChainS\_of\_inf\textquotesingle{}}
  (\href{https://github.com/jpbald93/paperl_clifford_chain/blob/d15464f0f4f4c8848e7d3e63dd68a03ce45433ae/Clifford/MorleyKernel.lean\#L4650}{\texttt{Clifford/\allowbreak MorleyKernel.lean}:4650}), which is the induction \texttt{cliffordChainS\_\allowbreak all\_\allowbreak of}
  (\href{https://github.com/jpbald93/paperl_clifford_chain/blob/d15464f0f4f4c8848e7d3e63dd68a03ce45433ae/Clifford/Mobius.lean\#L330}{\texttt{Clifford/\allowbreak Mobius.lean}:330}) instantiated at \texttt{E\ \allowbreak :=\ EptInf} with \texttt{h4\_\allowbreak pt\_\allowbreak inf}
  (\href{https://github.com/jpbald93/paperl_clifford_chain/blob/d15464f0f4f4c8848e7d3e63dd68a03ce45433ae/Clifford/MorleyKernel.lean\#L2566}{\texttt{Clifford/\allowbreak MorleyKernel.lean}:2566}) as its level-4 certificate; the \texttt{n\ =\ 3} base of that induction
  comes from \texttt{cliffordChain\_\allowbreak threeS} (\href{https://github.com/jpbald93/paperl_clifford_chain/blob/d15464f0f4f4c8848e7d3e63dd68a03ce45433ae/Clifford/Mobius.lean\#L181}{\texttt{Clifford/\allowbreak Mobius.lean}:181}). \textbf{No single declaration in the
  library states \texttt{∀\ n,\ CliffordChainS\ n}};
  the aggregate is their composition, which is stated and checked, at the pin, in
  \texttt{check/\allowbreak allN.\allowbreak lean} (\texttt{lake\ env\ lean} rc 0 in about 3 s, \textbf{std} axioms, 2026-11-24).
\item
  \textbf{C6 --- a uniform step.} The odd side and the even side are proved by separate inductions but
  from one algebraic engine: the Morley (interpolation) matrix \texttt{KP\ n\ c\ (conj\ c)} of size n$\times$(n+1),
  whose kernel vector at level n is the Clifford object. The odd step is the payload
  \texttt{hodd\_inf\_all\textquotesingle{}}; the even step reduces to a single scalar identity \texttt{EvenCornerQL} in a degenerate
  corner, which was proved in round 26 (267 added lines) and promoted into the library in round 26b.
\end{itemize}

A degeneracy that had survived three earlier rounds is worth naming because it is where the work
was: on the even side every branch closed except the corner \texttt{hmid\ =\ 0\ ∧\ V\_\allowbreak last\ =\ 0\ ∧\ P₀(d)\ ≠\ 0},
which forces the parent cross vector \texttt{V\ =\ 0} (from \texttt{hmid\ =\ 0} and \texttt{P₀(d)\ ≠\ 0} alone, by
\texttt{parent\_\allowbreak crossVec\_\allowbreak zero\_\allowbreak of\_\allowbreak mid}, \href{https://github.com/jpbald93/paperl_clifford_chain/blob/d15464f0f4f4c8848e7d3e63dd68a03ce45433ae/Clifford/MorleyKernel.lean\#L4998}{\texttt{Clifford/\allowbreak MorleyKernel.lean}:4998}; the packaging lemma
\texttt{heven\_\allowbreak inf\_\allowbreak of\_\allowbreak corner} additionally requires all grandchildren finite). The corner is not vacuous in
practice: a numerical search found configurations approximately in it, with the grandchildren's level
points collinear in the search (\texttt{1e-8} in the earlier, looser run, \texttt{3.\allowbreak 9e-5} for the tuple the JSON
stores); exact membership is not claimed from that search, and the theorem below does not need it. The recorded witness has \texttt{triP} 1e-2, \texttt{cycP} 1.3e-3, transversals $\ge$ 1e-2 and
all \texttt{\textbar{}c\textbar{}\ $\ge$\ 0.08}; an earlier, looser solver run had \texttt{triP} 0.14, \texttt{cycP} 0.024, transversals 0.04 (the
JSON records both runs separately; the stored tuple's \texttt{\textbar{}V\textbar{}} max is \texttt{4.\allowbreak 2e-10} and its subfamily
collinearity deviation \texttt{3.\allowbreak 9e-5}, while the \texttt{1e-9} residual and the \texttt{1e-8} collinearity figure belong
to the earlier run). That witness is \textbf{numerical evidence, not a certificate} --- small residuals do
not prove the exact vanishing of \texttt{V} --- and it is bundled at
\texttt{layer0/\allowbreak r25\_\allowbreak jack/\allowbreak RANK\_\allowbreak ROUTE\_\allowbreak COUNTEREXAMPLE.\allowbreak json} (see Appendix A). It argues against the tempting
rank route; it does not refute it, and nothing below rests on it. The exact statement proved later is
stronger than collinearity: \texttt{evenCorner\_\allowbreak of\_\allowbreak QL}
(\href{https://github.com/jpbald93/paperl_clifford_chain/blob/d15464f0f4f4c8848e7d3e63dd68a03ce45433ae/Clifford/MorleyKernel.lean\#L5287}{\texttt{Clifford/\allowbreak MorleyKernel.lean}:5287}) and \texttt{evenCornerQL\_\allowbreak proof} (\href{https://github.com/jpbald93/paperl_clifford_chain/blob/d15464f0f4f4c8848e7d3e63dd68a03ce45433ae/Clifford/MorleyKernel.lean\#L5539}{\texttt{Clifford/\allowbreak MorleyKernel.lean}:5539})
show that in the exact corner all \texttt{2k+5} grandchild level points \textbf{coincide} at the single point
\texttt{$-$w\_\{2k+4\}/w\_\{k+2\}}. Round 25 reduced the corner to the scalar identity \texttt{w\_\allowbreak last\ =\ 0} for the
subfamily cross vector \texttt{w}, and round 26 proved that identity by a polynomial argument on
dehomogenised coordinates: the row equations make \texttt{polyF\ :=\ $-$X$\cdot$Q$\cdot$\=P\^{}rev\ $-$\ P$\cdot$\=Q\^{}rev}, of degree
$\le$ 2k+5, vanish at 2k+6 pairwise non-proportional points, so it is identically zero and its top
coefficient \texttt{$-$W\_\{2k+5\}$\cdot$conj\ W\_0} vanishes. The mathematics of that argument is elementary; the
formalisation is not, and it is the one place where this development goes beyond a transcription of
the classical proofs.

\hypertarget{what-this-paper-does-not-claim}{%
\subsection{What this paper does not claim}\label{what-this-paper-does-not-claim}}

\begin{itemize}
\tightlist
\item
  \textbf{The at-infinity / improper-set classification is Carver's (1920), never ours.} Carver's \S{}5
  (pp.~160--161) proves the key-matrix criterion and its Theorem 13; \S{}6 (pp.~161--164) treats the
  limiting figures, and it is \S{}6 that states the inversion convention this paper's carrier is
  modelled on. The correspondence is not an identification: \S{}6 Type 2 explicitly admits a
  \textbf{point-circle at infinity}, and Types 3/3$'$/6 have partially or entirely indeterminate limiting
  figures, whereas \texttt{IsGCircleS} admits finite point-circles only and excludes \texttt{\{$\infty$\}}. Morley (1900)
  supplies the analytic method and the ``Clifford line'' phenomenon, Kantor (1878) the deltoid
  examples. In this paper these appear only as the \emph{reason} the sphere predicate is the correct one,
  and as the citation for the $\infty$ case. No classification theorem is formalised here: the Lean
  development proves the chain exists, including when the level object is a line or the point at
  infinity, but it does not state the characteristic-matrix criterion in general form. \textbf{OPEN}: a
  formal version of Carver's Theorem 13 is not part of this development.
\item
  \textbf{The mathematics is not new}: the all-n chain (Clifford 1870--71/1882; Morley 1900; Carver 1920;
  Longuet-Higgins 1972), the failure phenomena (Kantor 1878; Carver 1920), the M\"obius-plane repair
  (Carver 1920 \S{}6), the analytic method (Morley 1900) and the higher-polytope picture
  (Longuet-Higgins 1972) and the analytic, cross-ratio route (Bobenko and Schief,
  arXiv:1509.04109, published 2016) are all classical.
\item
  \textbf{First formalisation of Miquel is also not claimed}: atlas-lean, AFP 2026, Coq HSG and
  lean-genius all contain a Miquel theorem.
\item
  Our own earlier paper \textbf{Paper D} (Zenodo record 23653309) proves C$_4$ and C$_5$ in the \emph{finite}
  plane unconditionally; its C$_6$ statement carries a further non-degeneracy hypothesis (the
  non-collinearity of selected level-4 points in every five-subfamily), so it is not the same
  predicate as this paper's \texttt{CliffordChain\ 6}, which \S{}2.2 proves \textbf{false}. What is reused as base
  cases here is only the finite-plane chain for \textbf{n $\le$ 5}, transported by \texttt{chainS\_\allowbreak of\_\allowbreak chain}; the
  all-n $\ge$ 6 failure of the finite-plane predicate is \texttt{not\_\allowbreak cliffordChain\_\allowbreak of\_\allowbreak six\_\allowbreak le}
  (\href{https://github.com/jpbald93/paperl_clifford_chain/blob/d15464f0f4f4c8848e7d3e63dd68a03ce45433ae/Clifford/Classify.lean\#L219}{\texttt{Clifford/\allowbreak Classify.lean}:219}), which builds on Paper D's six-centre counterexample.
\end{itemize}


\hypertarget{statements}{%
\section{Statements}\label{statements}}

\hypertarget{the-carrier}{%
\subsection{The carrier}\label{the-carrier}}

Everything is stated for circles through the origin \texttt{M\ =\ 0} (a similarity moves any M there) with
centres \texttt{c\ \allowbreak :\ Fin\ n\ →\ ℂ}. The two ground-level notions are (\href{https://github.com/jpbald93/paperl_clifford_chain/blob/d15464f0f4f4c8848e7d3e63dd68a03ce45433ae/Clifford/Basic.lean\#L394}{\texttt{Clifford/\allowbreak Basic.lean}:394}, \href{https://github.com/jpbald93/paperl_clifford_chain/blob/d15464f0f4f4c8848e7d3e63dd68a03ce45433ae/Clifford/Basic.lean\#L398}{\texttt{Clifford/\allowbreak Basic.lean}:398})

\begin{Verbatim}
noncomputable def second (c d : ℂ) : ℂ := (conj c * d - c * conj d) / (conj c - conj d)

def Transversal (c d : ℂ) : Prop := conj c * d ≠ c * conj d
\end{Verbatim}

\texttt{second\ c\ d} is the second intersection of the two circles (besides \texttt{M\ =\ 0}), and \texttt{Transversal\ c\ d}
says the two circles meet transversally, equivalently that \texttt{M}, \texttt{c}, \texttt{d} are not collinear. General
position is \texttt{GenPosN} (\href{https://github.com/jpbald93/paperl_clifford_chain/blob/d15464f0f4f4c8848e7d3e63dd68a03ce45433ae/Clifford/Clean.lean\#L177}{\texttt{Clifford/\allowbreak Clean.lean}:177}):

\begin{Verbatim}
def GenPosN {n : ℕ} (c : Fin n → ℂ) : Prop :=
  (∀ i j, i ≠ j → Transversal (c i) (c j)) ∧
    ∀ i j k, i < j → j < k → GenTriple (c i) (c j) (c k)
\end{Verbatim}

with \texttt{GenTriple} (\href{https://github.com/jpbald93/paperl_clifford_chain/blob/d15464f0f4f4c8848e7d3e63dd68a03ce45433ae/Clifford/Basic.lean\#L402}{\texttt{Clifford/\allowbreak Basic.lean}:402})

\begin{Verbatim}
def GenTriple (a b c : ℂ) : Prop :=
  triP a b c (conj a) (conj b) (conj c) ≠ 0 ∧ cycP a b c (conj a) (conj b) (conj c) ≠ 0
\end{Verbatim}

Both conjuncts matter, and the second is not a formality: by \texttt{genTriple\_\allowbreak iff}
(\href{https://github.com/jpbald93/paperl_clifford_chain/blob/d15464f0f4f4c8848e7d3e63dd68a03ce45433ae/Clifford/Bridge.lean\#L110}{\texttt{Clifford/\allowbreak Bridge.lean}:110}) the pair says that the three centres are not collinear \textbf{and} that the
three second intersections are not collinear. The second conjunct is what guarantees that the
level-3 object is a genuine circle and never a line, which is why lines --- and the point $\infty$ --- can
first occur at levels 5 and 6, not at level 3.

The sphere carrier is defined in \texttt{Clifford/\allowbreak Mobius.\allowbreak lean}:

\begin{Verbatim}
def lineS (p v : ℂ) : Set (OnePoint ℂ) :=
  insert ∞ (((↑) : ℂ → OnePoint ℂ) '' {z | ∃ t : ℝ, z = p + t • v})

def IsGCircleS (X : Set (OnePoint ℂ)) : Prop :=
  (∃ o : ℂ, ∃ r : ℝ, X = ((↑) : ℂ → OnePoint ℂ) '' Metric.sphere o r) ∨
    ∃ p v : ℂ, v ≠ 0 ∧ X = lineS p v
\end{Verbatim}

so the first disjunct is \emph{any} \texttt{Metric.\allowbreak sphere\ o\ r} with \texttt{r} an arbitrary real: \texttt{r\ \textgreater{}\ 0} a genuine
circle, \texttt{r\ =\ 0} a single point (a point-circle, which the classical convention admits), \texttt{r\ \textless{}\ 0} the
empty set --- which the statements never produce, since an empty set cannot contain a level point. The
second disjunct is a line of \texttt{ℂ} together with \texttt{∞}. The development carries the empty case rather
than excluding it, so the prose matches the predicate and not an idealised version of it. The chain recursion is

\begin{Verbatim}
def ChainS : (n : ℕ) → (Fin n → ℂ) → Set (OnePoint ℂ) → Prop
  | 0, _, X => X = {((0 : ℂ) : OnePoint ℂ)}
  | 1, c, X => X = ((↑) : ℂ → OnePoint ℂ) '' Metric.sphere (c 0) ‖c 0‖
  | 2, c, X => X = {((second (c 0) (c 1) : ℂ) : OnePoint ℂ)}
  | n + 3, c, X =>
    if n % 2 = 0 then
      IsGCircleS X ∧ ∀ i : Fin (n + 3), ∃ p, ChainS (n + 2) (c ∘ i.succAbove) {p} ∧ p ∈ X
    else
      ∃ p, X = {p} ∧ ∀ i : Fin (n + 3), ∃ Y, ChainS (n + 2) (c ∘ i.succAbove) Y ∧ p ∈ Y
\end{Verbatim}

and the theorem is

\begin{Verbatim}
def CliffordChainS (n : ℕ) : Prop := ∀ c : Fin n → ℂ, GenPosN c → ∃ X, ChainS n c X
\end{Verbatim}

\texttt{ChainS} differs from the finite-plane \texttt{Chain} (\href{https://github.com/jpbald93/paperl_clifford_chain/blob/d15464f0f4f4c8848e7d3e63dd68a03ce45433ae/Clifford/Clean.lean\#L165}{\texttt{Clifford/\allowbreak Clean.lean}:165}) in two respects: its
carrier is \texttt{Set\ (OnePoint\ ℂ)} throughout, so the level-2 point is a point of the sphere; and at odd
levels the object must be a generalized circle \emph{on the sphere} (\texttt{IsGCircleS}, lines containing \texttt{∞}
allowed), while even levels quantify over sphere points. The combined effect is that nothing forces a
later object to be finite.
\texttt{liftS\ Y} (\href{https://github.com/jpbald93/paperl_clifford_chain/blob/d15464f0f4f4c8848e7d3e63dd68a03ce45433ae/Clifford/Mobius.lean\#L51}{\texttt{Clifford/\allowbreak Mobius.lean}:51}) is the closure of a generalized circle of \texttt{ℂ} in the sphere ---
the set with \texttt{∞} adjoined iff \texttt{Y} is unbounded --- and \texttt{chainS\_\allowbreak of\_\allowbreak chain}
(\href{https://github.com/jpbald93/paperl_clifford_chain/blob/d15464f0f4f4c8848e7d3e63dd68a03ce45433ae/Clifford/Mobius.lean\#L119}{\texttt{Clifford/\allowbreak Mobius.lean}:119}) transports any finite chain object along it.

The level point is read off the kernel vector of a Morley matrix. \texttt{crossVec}
(\href{https://github.com/jpbald93/paperl_clifford_chain/blob/d15464f0f4f4c8848e7d3e63dd68a03ce45433ae/Clifford/MorleyKernel.lean\#L159}{\texttt{Clifford/\allowbreak MorleyKernel.lean}:159}) is the vector of signed maximal minors of an n$\times$(n+1) matrix,

\begin{Verbatim}
def crossVec {R : Type*} [CommRing R] {n : ℕ} (A : Matrix (Fin n) (Fin (n + 1)) R) :
    Fin (n + 1) → R :=
  fun j => (-1) ^ (j : ℕ) * (A.submatrix id j.succAbove).det
\end{Verbatim}

\texttt{KP\ n\ x\ y} (\href{https://github.com/jpbald93/paperl_clifford_chain/blob/d15464f0f4f4c8848e7d3e63dd68a03ce45433ae/Clifford/MorleyKernel.lean\#L214}{\texttt{Clifford/\allowbreak MorleyKernel.lean}:214}) is the row-scaled n$\times$(n+1) Morley (interpolation)
matrix in \texttt{x\ =\ c}, \texttt{y\ =\ conj\ c}, with entries

\begin{Verbatim}
def KP {R : Type*} [CommRing R] (n : ℕ) (x y : Fin n → R) : Matrix (Fin n) (Fin (n + 1)) R :=
  fun k j => if (j : ℕ) ≤ n / 2 then (-x k) ^ (j : ℕ) * y k ^ (n - n / 2 - j)
    else -((-x k) ^ ((j : ℕ) - (n / 2 + 1)) * y k ^ (n - n / 2 - 1 - ((j : ℕ) - (n / 2 + 1))))
\end{Verbatim}

It is row \texttt{k} of the polynomial matrix \texttt{Kf\ n\ (aOf\ c)\ (tOf\ c)} (\href{https://github.com/jpbald93/paperl_clifford_chain/blob/d15464f0f4f4c8848e7d3e63dd68a03ce45433ae/Clifford/MorleyKernel.lean\#L95}{\texttt{Clifford/\allowbreak MorleyKernel.lean}:95}, with
\texttt{aOf\ c\ =\ 1\ /\allowbreak \ conj\ c} and \texttt{tOf\ c\ =\ −c\ /\allowbreak \ conj\ c}, \href{https://github.com/jpbald93/paperl_clifford_chain/blob/d15464f0f4f4c8848e7d3e63dd68a03ce45433ae/Clifford/MorleyKernel.lean\#L219}{\texttt{Clifford/\allowbreak MorleyKernel.lean}:219}) multiplied by
\texttt{(conj\ c\ k)\^{}(n\ $-$\ n/2)}. For \textbf{nonzero centres} row scaling does not change the kernel
(\texttt{KP\_\allowbreak mulVec\_\allowbreak eq\_\allowbreak zero\_\allowbreak iff}, \href{https://github.com/jpbald93/paperl_clifford_chain/blob/d15464f0f4f4c8848e7d3e63dd68a03ce45433ae/Clifford/MorleyKernel.lean\#L236}{\texttt{Clifford/\allowbreak MorleyKernel.lean}:236}); nonzero centres follow from \texttt{GenPosN} at
every level where the readout is used (\texttt{hc\_\allowbreak of\_\allowbreak genPos}, \href{https://github.com/jpbald93/paperl_clifford_chain/blob/d15464f0f4f4c8848e7d3e63dd68a03ce45433ae/Clifford/MorleyKernel.lean\#L2571}{\texttt{Clifford/\allowbreak MorleyKernel.lean}:2571}) but not at
\texttt{n\ =\ 1}, where \texttt{c\ =\ 0} satisfies \texttt{GenPosN} and \texttt{KP\ =\ {[}0,\ $-$1{]}} while \texttt{Kf\ =\ {[}1,\ 0{]}}, so the kernel
equivalence is stated with that hypothesis. The readouts are
(\href{https://github.com/jpbald93/paperl_clifford_chain/blob/d15464f0f4f4c8848e7d3e63dd68a03ce45433ae/Clifford/MorleyKernel.lean\#L452}{\texttt{Clifford/\allowbreak MorleyKernel.lean}:452}, \href{https://github.com/jpbald93/paperl_clifford_chain/blob/d15464f0f4f4c8848e7d3e63dd68a03ce45433ae/Clifford/MorleyKernel.lean\#L2549}{\texttt{Clifford/\allowbreak MorleyKernel.lean}:2549})

\begin{Verbatim}
noncomputable def levelPointN (n : ℕ) (v : Fin (n + 1) → ℂ) : ℂ :=
  -v (Fin.last n) / v ⟨n / 2, by omega⟩

noncomputable def EptInf (n : ℕ) (c : Fin n → ℂ) : OnePoint ℂ :=
  if crossVec (KP n c (fun j => conj (c j))) ⟨n / 2, by omega⟩ = 0 then OnePoint.infty
  else ((levelPointN n (crossVec (KP n c (fun j => conj (c j)))) : ℂ) : OnePoint ℂ)
\end{Verbatim}

So \texttt{EptInf\ n\ c} is the projective point \texttt{{[}$-$v\_last\ :\ v\_\{n/2\}{]}} of \texttt{v\ =\ crossVec\ (KP\ n\ c\ \=c)}: a
vanishing middle coordinate is sent to \texttt{∞}. This includes the case of the \textbf{whole vector zero},
which is not a projective coordinate at all; sending it to \texttt{∞} is the development's chosen
totalization, not a consequence of projective geometry. That branch is not vacuous: for the parabola
configuration of \S{}2.4 (centres \texttt{2t\ −\ 2it²}, \texttt{t\ =\ 1,\ …,\ 6}) the middle coordinate of the level-6
parent cross vector vanishes and, by \texttt{parent\_\allowbreak crossVec\_\allowbreak zero\_\allowbreak of\_\allowbreak mid}, so does the whole vector --- two
exact evaluations at the pinned tree, recorded in \texttt{REVIEW\_\allowbreak TRIAGE\_\allowbreak v03b.\allowbreak md} (\textbf{std} axioms); general
position for that tuple is proved in \texttt{check/\allowbreak not\_\allowbreak UniqEven.\allowbreak lean}. \texttt{Counter6} itself is \emph{not} an instance: its parent vector has last coordinate
\texttt{−238878720000}, checked exactly. The choice of totalization is therefore visible in the statements
rather than hidden.

\hypertarget{why-the-sphere-is-needed-the-finite-plane-is-false}{%
\subsection{Why the sphere is needed: the finite plane is false}\label{why-the-sphere-is-needed-the-finite-plane-is-false}}

The finite-plane predicate is

\begin{Verbatim}
def CliffordChain (n : ℕ) : Prop := ∀ c : Fin n → ℂ, GenPosN c → ∃ X, Chain n c X
\end{Verbatim}

(\href{https://github.com/jpbald93/paperl_clifford_chain/blob/d15464f0f4f4c8848e7d3e63dd68a03ce45433ae/Clifford/Clean.lean\#L184}{\texttt{Clifford/\allowbreak Clean.lean}:184}), and the development proves

\begin{Verbatim}
theorem cliffordChain_iff (n : ℕ) : CliffordChain n ↔ n ≤ 5        -- Clifford/Classify.lean:227
\end{Verbatim}

by exhibiting six centres in general position with no level-6 object:
\texttt{Counter6.\allowbreak cx\ =\ ⟨3,−6⟩,\ ⟨6,−3⟩,\ ⟨1,2⟩,\ ⟨−2,−1⟩,\ ⟨0,−4⟩,\ ⟨4,0⟩} (exact Gaussian integers,
\href{https://github.com/jpbald93/paperl_clifford_chain/blob/d15464f0f4f4c8848e7d3e63dd68a03ce45433ae/Clifford/Counter6.lean\#L19}{\texttt{Clifford/\allowbreak Counter6.lean}:19}). For them every level-4 point is forced to be the Miquel point; as an
explanation rather than a quoted theorem, the six level-5 objects are lines (five distinct collinear
Miquel points each --- the lines themselves are \texttt{line0}\ldots{}\texttt{line5}, \href{https://github.com/jpbald93/paperl_clifford_chain/blob/d15464f0f4f4c8848e7d3e63dd68a03ce45433ae/Clifford/Mobius.lean\#L492}{\texttt{Clifford/\allowbreak Mobius.lean}:492} through
\href{https://github.com/jpbald93/paperl_clifford_chain/blob/d15464f0f4f4c8848e7d3e63dd68a03ce45433ae/Clifford/Mobius.lean\#L627}{\texttt{Clifford/\allowbreak Mobius.lean}:627}), and three of those lines are not concurrent (\texttt{Counter6.\allowbreak not\_\allowbreak chain6},
\href{https://github.com/jpbald93/paperl_clifford_chain/blob/d15464f0f4f4c8848e7d3e63dd68a03ce45433ae/Clifford/Counter6.lean\#L182}{\texttt{Clifford/\allowbreak Counter6.lean}:182}; the universal form \texttt{Counter6.\allowbreak not\_\allowbreak cliffordChain\_\allowbreak six}, \href{https://github.com/jpbald93/paperl_clifford_chain/blob/d15464f0f4f4c8848e7d3e63dd68a03ce45433ae/Clifford/Counter6.lean\#L227}{\texttt{Clifford/\allowbreak Counter6.lean}:227}). What
the proof of \texttt{not\_\allowbreak chain6} actually uses is the contradiction from the \texttt{cD\ =\ 0} relations at three of
the six deletions (\texttt{i\ =\ 4,\ 5,\ 0}). This is Paper D's failure phenomenon, and on the sphere it
disappears rather than contradicting anything: the level-6 object is \texttt{\{$\infty$\}},

\begin{Verbatim}
theorem chainS_six : ChainS 6 Counter6.cx {∞}                 -- Clifford/Mobius.lean:651
theorem chainS_six_unique {X : Set (OnePoint ℂ)} (h : ChainS 6 Counter6.cx X) : X = {∞}
theorem chainS_six_iff (X : Set (OnePoint ℂ)) : ChainS 6 Counter6.cx X ↔ X = {∞}
\end{Verbatim}

The pair \texttt{cliffordChain\_\allowbreak iff} / \texttt{chainS\_\allowbreak six\_\allowbreak unique} is the precise sense in which the old predicate
was too small: the sphere adds the level-6 point this configuration was missing. It is \emph{inspired by}
Carver's reading (``the plane has only one point at infinity \ldots{} circle includes straight lines and
point-circles'', \texttt{SCOUT.\allowbreak md} \S{}0 row 18) but not identical to it --- see \S{}1.2 --- and the uniqueness of
\texttt{\{$\infty$\}} here is a property of this one configuration, not of the predicate (\S{}5.4).

\hypertarget{the-odd-side-is-hypothesis-free}{%
\subsection{The odd side is hypothesis-free}\label{the-odd-side-is-hypothesis-free}}

For odd levels the payload is an existence statement for a generalized circle through the canonical
level-(n$-$1) points, and it holds for every \texttt{GenPosN} configuration with no side conditions
(\href{https://github.com/jpbald93/paperl_clifford_chain/blob/d15464f0f4f4c8848e7d3e63dd68a03ce45433ae/Clifford/MorleyKernel.lean\#L4643}{\texttt{Clifford/\allowbreak MorleyKernel.lean}:4643}):

\begin{Verbatim}
theorem hodd_inf_all' :
    ∀ (k : ℕ) (d : Fin (2 * k + 5) → ℂ), GenPosN d →
      ∃ Y, Clifford.IsGCircleS Y ∧
        ∀ i : Fin (2 * k + 5), EptInf (2 * k + 4) (d ∘ i.succAbove) ∈ Y
\end{Verbatim}

It is obtained from the three-case split \texttt{hodd\_\allowbreak inf\_\allowbreak P0\_\allowbreak eq} / \texttt{hodd\_\allowbreak inf\_\allowbreak P0\_\allowbreak ne} /
\texttt{hodd\_\allowbreak inf\_\allowbreak P0\_\allowbreak ne\_\allowbreak some} on \texttt{P₀}. Note what the statement \emph{says}: it names the canonical point
\texttt{EptInf\ (2k+4)\ (d\ ∘\ i.\allowbreak succAbove)} rather than asserting that some point of that level lies on \texttt{Y}.
What turns these named readouts into actual chain objects is a separate step: \texttt{h4\_\allowbreak pt\_\allowbreak inf}
(\href{https://github.com/jpbald93/paperl_clifford_chain/blob/d15464f0f4f4c8848e7d3e63dd68a03ce45433ae/Clifford/MorleyKernel.lean\#L2566}{\texttt{Clifford/\allowbreak MorleyKernel.lean}:2566}) certifies the level-4 case (\texttt{ChainS\ 4\ c\ \{EptInf\ 4\ c\}} under
\texttt{GenPosN}), and \texttt{cliffordChainS\_\allowbreak all\_\allowbreak of} (\href{https://github.com/jpbald93/paperl_clifford_chain/blob/d15464f0f4f4c8848e7d3e63dd68a03ce45433ae/Clifford/Mobius.lean\#L330}{\texttt{Clifford/\allowbreak Mobius.lean}:330}) is the two-step induction that
rebuilds \texttt{ChainS\ (2k+4)\ c\ \{EptInf\ (2k+4)\ c\}} from that base case plus the odd and even payloads. A
reader who wants the chain itself must take the three together; \S{}5.3 records exactly what the
layer-13 comparison did with the payloads. That the odd side needs no hypothesis beyond general
position is a statement about the development, not about the mathematics: Carver's improper sets
exist on the odd side too, and the reason we can state the odd side cleanly is that \texttt{IsGCircleS}
admits the line as a legitimate level object.

\hypertarget{the-even-side-and-the-corner}{%
\subsection{The even side, and the corner}\label{the-even-side-and-the-corner}}

On the even side the deliverable is: for every configuration of \texttt{2k+6} circles and every choice \texttt{i}
of parent object --- the subfamily obtained by deleting \texttt{c\ i}, whose object sits at level \texttt{2k+5}
(level 5 only when \texttt{k\ =\ 0}) --- there is a generalized circle \texttt{Y} through the canonical level-(2k+6)
point which contains all \texttt{2k+5} canonical level-(2k+4) points of the sub-configurations. Written out
(\href{https://github.com/jpbald93/paperl_clifford_chain/blob/d15464f0f4f4c8848e7d3e63dd68a03ce45433ae/Clifford/MorleyKernel.lean\#L5586}{\texttt{Clifford/\allowbreak MorleyKernel.lean}:5586}):

\begin{Verbatim}
theorem heven_inf_all :
    ∀ (k : ℕ) (c : Fin (2*k+6) → ℂ), GenPosN c → ∀ i : Fin (2*k+6),
      ∃ Y, Clifford.IsGCircleS Y ∧ EptInf (2*k+6) c ∈ Y ∧
        ∀ j : Fin (2*k+5), EptInf (2*k+4) ((c ∘ i.succAbove) ∘ j.succAbove) ∈ Y
\end{Verbatim}

As in \S{}2.3, the theorem names the canonical points rather than saying ``some point of the level'', and
again the named readouts are certified as chain objects through \texttt{h4\_\allowbreak pt\_\allowbreak inf} and the induction --- not
by the payload alone. The proof splits four ways; three branches are unconditional: \texttt{P₀(d)\ =\ 0}
(\texttt{heven\_\allowbreak line\_\allowbreak of\_\allowbreak P0}), \texttt{hmid\ ≠\ 0} with all children finite (\texttt{heven\_\allowbreak kernel}), and some child
\texttt{P\_\{k+2\}\ =\ 0} (\texttt{children\_\allowbreak eq\_\allowbreak of\_\allowbreak P3\_\allowbreak zero}: there the \textbf{finite} child points coincide; a child with a
vanishing middle coordinate is \texttt{∞} by \texttt{EptInf}, and \texttt{heven\_\allowbreak of\_\allowbreak common} builds a line through the common
finite point, any infinite children and the parent --- no non-uniqueness is inferred in this branch). The fourth is the corner. It is named twice, first as the geometric statement
(\href{https://github.com/jpbald93/paperl_clifford_chain/blob/d15464f0f4f4c8848e7d3e63dd68a03ce45433ae/Clifford/MorleyKernel.lean\#L5253}{\texttt{Clifford/\allowbreak MorleyKernel.lean}:5253})

\begin{Verbatim}
def EvenCorner : Prop :=
  ∀ (k : ℕ) (c : Fin (2*k+6) → ℂ), GenPosN c → ∀ i : Fin (2*k+6),
    crossVec (KP (2*k+6) c (fun j => conj (c j))) = 0 →
    crossVec (KP (2*k+5) (c ∘ i.succAbove) (fun r => conj ((c ∘ i.succAbove) r))) 0 ≠ 0 →
    (∀ j : Fin (2*k+5), crossVec (KP (2*k+4) ((c ∘ i.succAbove) ∘ j.succAbove)
        (fun r => conj (((c ∘ i.succAbove) ∘ j.succAbove) r))) ⟨k+2, by omega⟩ ≠ 0) →
    ∃ p u : ℂ, u ≠ 0 ∧ ∀ j : Fin (2*k+5), ∃ t : ℝ,
      levelPointN (2*k+4) (crossVec (KP (2*k+4) ((c ∘ i.succAbove) ∘ j.succAbove)
        (fun r => conj (((c ∘ i.succAbove) ∘ j.succAbove) r)))) = p + (t : ℂ) * u
\end{Verbatim}

and then, after round 25, as the scalar identity that implies it (\href{https://github.com/jpbald93/paperl_clifford_chain/blob/d15464f0f4f4c8848e7d3e63dd68a03ce45433ae/Clifford/MorleyKernel.lean\#L5317}{\texttt{Clifford/\allowbreak MorleyKernel.lean}:5317}):

\begin{Verbatim}
def EvenCornerQL : Prop := … (same hypotheses) …
    crossVec (KP (2*k+5) (c ∘ i.succAbove) (fun r => conj ((c ∘ i.succAbove) r))) (Fin.last _) = 0

theorem evenCorner_of_QL …                                     -- Clifford/MorleyKernel.lean:5287
theorem evenCorner_of_cornerQL (hQ : EvenCornerQL) : EvenCorner  -- Clifford/MorleyKernel.lean:5325
\end{Verbatim}

One property of \texttt{EvenCornerQL} is worth stating because it looks like a gap and is not one: the
grandchild-finiteness hypothesis (the third premise, \texttt{hall}) is never used by \texttt{evenCornerQL\_\allowbreak proof}
(\href{https://github.com/jpbald93/paperl_clifford_chain/blob/d15464f0f4f4c8848e7d3e63dd68a03ce45433ae/Clifford/MorleyKernel.lean\#L5539}{\texttt{Clifford/\allowbreak MorleyKernel.lean}:5539} introduces it as \texttt{\_\allowbreak hall}), so the identity holds without it. Layer
6 could not flag it, because \texttt{EvenCornerQL} bundles its hypotheses inside a \texttt{Prop} definition rather
than as top-level binders, and the ablation pass only sees binders. \texttt{hall} \emph{is} used downstream, by
\texttt{evenCorner\_\allowbreak of\_\allowbreak QL}, to deduce \texttt{EvenCorner}.

Round 26 proved \texttt{EvenCornerQL} itself (\texttt{evenCornerQL\_\allowbreak proof}, \href{https://github.com/jpbald93/paperl_clifford_chain/blob/d15464f0f4f4c8848e7d3e63dd68a03ce45433ae/Clifford/MorleyKernel.lean\#L5539}{\texttt{Clifford/\allowbreak MorleyKernel.lean}:5539}, +267 lines over the round-25
prefix, which is byte-identical on its first 5344 lines \textbf{in the scratch file}
\texttt{scratch/\allowbreak ChainEvenCornerQL.\allowbreak lean}; the promoted file's lines 1006--2005 differ from it by the gate edits
of \S{}4), and the chain follows with no hypotheses left:

\begin{Verbatim}
theorem EvenCorner_proof : EvenCorner                       -- Clifford/MorleyKernel.lean:5584
theorem heven_inf_all : …                                   -- Clifford/MorleyKernel.lean:5586
theorem cliffordChainS_inf (n : ℕ) (hn : 3 ≤ n) : CliffordChainS n   -- Clifford/MorleyKernel.lean:5592
\end{Verbatim}

\textbf{The even-level point is not unique, and no uniqueness theorem is available.} \texttt{EptInf} selects one
canonical point per level, but \texttt{ChainS} is existential at every level, and in degenerate
configurations it admits more than one even-level object. The counterexample checked for this draft
(\texttt{check/\allowbreak not\_\allowbreak UniqEven.\allowbreak lean}, \texttt{lake\ env\ lean} rc 0 in about 28 s, \textbf{std} axioms, 2026-11-24) takes the six
centres \texttt{2t\ −\ 2it²}, \texttt{t\ =\ 1.\allowbreak .\allowbreak 6}: \texttt{GenPosN} holds, all fifteen four-subfamily Miquel points are \texttt{4i},
every five-subfamily admits the line \texttt{lineS\ (4i)\ 1}, and that line contains both \texttt{4i} and \texttt{∞}, so
both \texttt{\{4i\}} and \texttt{\{$\infty$\}} satisfy \texttt{ChainS\ 6}. Hence

\begin{Verbatim}
theorem not_UniqEven : ¬ (∀ (k : ℕ) (d : Fin (2*k+4) → ℂ), GenPosN d → ∀ p : OnePoint ℂ,
    ChainS (2*k+4) d {p} → p = EptInf (2*k+4) d)
\end{Verbatim}

A uniqueness statement of that form is false under the current definitions; any future one needs
extra hypotheses or a different selection predicate. What the development supplies is a choice of
canonical points (\texttt{EptInf}) together with the existential chain objects --- not a unique classical or
limiting object in degenerate cases. \S{}5.3 and \S{}5.4 record what this does and does not settle about
the prose.

\hypertarget{the-theorem}{%
\subsection{The theorem}\label{the-theorem}}

\begin{Verbatim}
theorem cliffordChainS_inf (n : ℕ) (hn : 3 ≤ n) : CliffordChainS n
\end{Verbatim}

with the levels \texttt{n\ ≤\ 5} supplied by \texttt{cliffordChainS\_\allowbreak of\_\allowbreak le\_\allowbreak five} (\href{https://github.com/jpbald93/paperl_clifford_chain/blob/d15464f0f4f4c8848e7d3e63dd68a03ce45433ae/Clifford/Mobius.lean\#L178}{\texttt{Clifford/\allowbreak Mobius.lean}:178}), which
derives them from the finite-plane theorem through \texttt{cliffordChainS\_\allowbreak of\_\allowbreak cliffordChain}
(\href{https://github.com/jpbald93/paperl_clifford_chain/blob/d15464f0f4f4c8848e7d3e63dd68a03ce45433ae/Clifford/Mobius.lean\#L157}{\texttt{Clifford/\allowbreak Mobius.lean}:157}) and the transport \texttt{chainS\_\allowbreak of\_\allowbreak chain} (\href{https://github.com/jpbald93/paperl_clifford_chain/blob/d15464f0f4f4c8848e7d3e63dd68a03ce45433ae/Clifford/Mobius.lean\#L119}{\texttt{Clifford/\allowbreak Mobius.lean}:119});
\texttt{chainS\_\allowbreak of\_\allowbreak chain} and its singleton form \texttt{chainS\_\allowbreak of\_\allowbreak chain\_\allowbreak pt} (\href{https://github.com/jpbald93/paperl_clifford_chain/blob/d15464f0f4f4c8848e7d3e63dd68a03ce45433ae/Clifford/Mobius.lean\#L152}{\texttt{Clifford/\allowbreak Mobius.lean}:152}) are the
two general transport implications, and it is \texttt{cliffordChainS\_\allowbreak of\_\allowbreak le\_\allowbreak five} --- not the transports ---
that gives the sphere base cases. Inside \texttt{cliffordChainS\_\allowbreak inf}, \texttt{n\ =\ 3} uses \texttt{cliffordChain\_\allowbreak threeS}
(\href{https://github.com/jpbald93/paperl_clifford_chain/blob/d15464f0f4f4c8848e7d3e63dd68a03ce45433ae/Clifford/Mobius.lean\#L181}{\texttt{Clifford/\allowbreak Mobius.lean}:181}) and \texttt{n\ =\ 4} uses \texttt{h4\_\allowbreak pt\_\allowbreak inf}, both transported. The aggregate claim for
every \texttt{n\ ≥\ 1} is the composition of the two, \texttt{check/\allowbreak allN.\allowbreak lean} (\S{}1.1).

Read out in prose, for every \texttt{n\ ≥\ 1} and every \texttt{GenPosN} family of n circles through M: a Clifford
object of level n exists on the Riemann sphere; for odd \texttt{n\ ≥\ 3} it is a generalized circle
containing, for each \texttt{i}, some level-(n$-$1) object of the family with \texttt{c\ i} deleted; for even
\texttt{n\ ≥\ 4} it is a point (possibly $\infty$) lying on, for each \texttt{i}, some level-(n$-$1) object of that family.
The indefinite article is forced by \texttt{ChainS}: the recursion quantifies existentially at every level,
so neither the lower-level objects nor the even-level point are determined in general (\S{}5.4). All
four terminal theorems --- \texttt{evenCornerQL\_\allowbreak proof}, \texttt{EvenCorner\_\allowbreak proof}, \texttt{heven\_\allowbreak inf\_\allowbreak all},
\texttt{cliffordChainS\_\allowbreak inf} --- depend on \textbf{std} only, as printed by \texttt{\#print\ axioms} at the end of the file
and re-checked by the gate.


\hypertarget{the-formalisation-map}{%
\section{The formalisation map}\label{the-formalisation-map}}

The development is a git repository, published at
\href{https://github.com/jpbald93/paperl_clifford_chain}{github.com/jpbald93/paperl\_clifford\_chain}, with
the audited tree at commit \texttt{d15464f} (156 files). The per-file sha256 table below is the identity of
that tree and the table every citation in this draft was checked against; the eleven hashes were
re-checked against a fresh clone of the published repository on 2026-11-26, immediately after the
push. Toolchain:
\texttt{leanprover/\allowbreak lean4\allowbreak :v4.\allowbreak 33.\allowbreak 1} (\texttt{lean-toolchain}); Mathlib pin
\texttt{0df444a360eaa60ab8c11dca51a86af692955474} (\texttt{lake-manifest.\allowbreak json}). File sizes and hashes of the ten
source files plus the gate:

{\scriptsize\begin{longtable}[]{@{}
  >{\raggedright\arraybackslash}p{(\columnwidth - 4\tabcolsep) * \real{0.3333}}
  >{\raggedright\arraybackslash}p{(\columnwidth - 4\tabcolsep) * \real{0.3333}}
  >{\raggedright\arraybackslash}p{(\columnwidth - 4\tabcolsep) * \real{0.3333}}@{}}
\toprule\noalign{}
\begin{minipage}[b]{\linewidth}\raggedright
file
\end{minipage} & \begin{minipage}[b]{\linewidth}\raggedright
lines
\end{minipage} & \begin{minipage}[b]{\linewidth}\raggedright
sha256
\end{minipage} \\
\midrule\noalign{}
\endhead
\bottomrule\noalign{}
\endlastfoot
\texttt{Clifford/\allowbreak Basic.\allowbreak lean} & 724 & \texttt{d5c14edb\allowbreak f2da8ec8\allowbreak 288484b4\allowbreak eeee54e3\allowbreak 0e85cf04\allowbreak e5c440c4\allowbreak 06b4c714\allowbreak 4ca0cb16} \\
\texttt{Clifford/\allowbreak General.\allowbreak lean} & 99 & \texttt{90b42ad2\allowbreak f407f49b\allowbreak 7ebc9555\allowbreak 73af3823\allowbreak 727cef54\allowbreak 03047ccd\allowbreak 2467222a\allowbreak fd36eeff} \\
\texttt{Clifford/\allowbreak Bridge.\allowbreak lean} & 121 & \texttt{93eb80b6\allowbreak 39791dc7\allowbreak e271249a\allowbreak 79e0439d\allowbreak d3fd08c8\allowbreak 31eb71e3\allowbreak 67129c52\allowbreak c8cecf5b} \\
\texttt{Clifford/\allowbreak Clean.\allowbreak lean} & 257 & \texttt{2a7cf40b\allowbreak 4c473e56\allowbreak c97db8b1\allowbreak 485aa5ad\allowbreak 468c9cbb\allowbreak 5c9a56b5\allowbreak b7ac9ada\allowbreak f33cc463} \\
\texttt{Clifford/\allowbreak Mobius.\allowbreak lean} & 787 & \texttt{f8c9f013\allowbreak 188f7767\allowbreak b8c2f585\allowbreak 0b0f761c\allowbreak fb689d0e\allowbreak 490a782d\allowbreak 5c9e10dd\allowbreak aec83200} \\
\texttt{Clifford/\allowbreak Counter6.\allowbreak lean} & 233 & \texttt{d6bf1cc9\allowbreak 13841fd4\allowbreak 87c539c6\allowbreak 1e4db620\allowbreak 54e62d4b\allowbreak 861af3a3\allowbreak cb9eaac8\allowbreak cf35d57a} \\
\texttt{Clifford/\allowbreak Classify.\allowbreak lean} & 239 & \texttt{ba4fc1b3\allowbreak c567f5d6\allowbreak 9ca33ea2\allowbreak bcbab0df\allowbreak aed3a54a\allowbreak 779e3405\allowbreak efdb4241\allowbreak be65b83b} \\
\texttt{Clifford/\allowbreak MorleyKernel.\allowbreak lean} & 5600 & \texttt{dce7b011\allowbreak a8f22469\allowbreak 9f7ea1e8\allowbreak 79a159fc\allowbreak 8066e15f\allowbreak 3cb8c6e7\allowbreak 664c03a7\allowbreak 79990a76} \\
\texttt{Clifford/\allowbreak BridgeL10.\allowbreak lean} & 251 & \texttt{097722f3\allowbreak d2ff9949\allowbreak 65dcb1cc\allowbreak ff5180d8\allowbreak a62404ea\allowbreak 3dfc362b\allowbreak ef345934\allowbreak c1ab6054} \\
\texttt{Clifford.\allowbreak lean} (root) & --- & \texttt{8bb0f5f0\allowbreak d4663cb4\allowbreak e026980d\allowbreak 88fae6e2\allowbreak 4d0c2046\allowbreak b32062fb\allowbreak 6cce25a9\allowbreak e465f046} \\
\texttt{gate.\allowbreak sh} & --- & \texttt{2633ec93\allowbreak 10d126ba\allowbreak e9c1ff14\allowbreak b9ede499\allowbreak e7ffe85f\allowbreak c5749e0d\allowbreak 2a1de6cc\allowbreak 09881c89} \\
\end{longtable}}

Total: 8,311 lines of Lean across the nine \texttt{Clifford/\allowbreak *.\allowbreak lean} files, of which the layer-10 bridge file
contributes 251 (\texttt{MorleyKernel.\allowbreak lean} alone is 5,600 lines; before promotion the scratch file was 5,611 --- a
5,344-line prefix plus the 267-line round-26 proof --- and the two gate-forced edits of \S{}4 took it to
the 5,600 pinned here).

Every \texttt{file\allowbreak :line} citation in this draft --- prose, code comments, comma lists and tables alike --- is
machine-checked by \texttt{namegate.\allowbreak py} (layer 11, \texttt{L11\_\allowbreak NAMEGATE.\allowbreak md}): \textbf{143 occurrences of 98 distinct
\texttt{Clifford/\allowbreak File.\allowbreak lean\allowbreak :N} pairs}, each resolving to a declaration in the pinned tree, and the 25 checks
written in the form \texttt{\textasciigrave{}name\textasciigrave{}\ (\textasciigrave{}file:line\textasciigrave{})} additionally checked against the name the draft
attributes to them. Bare single-line citations have been removed; the eight ``lines A--B'' range citations left (four distinct spans)
name a region of a file or of a log, not a declaration. The script counts a citation inside a table
row or a code comment like any other, which is what the earlier 66-citation figure did not do (it
read only backticked citations in prose). What the gate checks is \textbf{names and lines, not
signatures}: \S{}3.3 anchors the carrier and statement declarations, \S{}2.1 quotes their definitions and
\S{}2.3--2.4 their statements, and the hypothesis-level audit is layer 12's (\texttt{L12\_\allowbreak REVIEW.\allowbreak md}, which
reviewed v0.1 --- see \S{}5.11).

\textbf{Availability.} The repository holds the ten Lean source files, \texttt{gate.\allowbreak sh}, \texttt{tests/\allowbreak tamper.\allowbreak sh},
\texttt{lakefile.\allowbreak toml}, \texttt{lean-toolchain} and \texttt{lake-manifest.\allowbreak json} at the pinned commit, and the two check
files of \S{}1.1 and \S{}5.4 in \texttt{check/\allowbreak } on \texttt{main}. \texttt{lake\ exe\ cache\ get} fetches the prebuilt Mathlib,
\texttt{lake\ build} compiles the library (the manifest, not the lakefile's version tag, fixes Mathlib at
\texttt{0df444a3…}), \texttt{.\allowbreak /\allowbreak gate.\allowbreak sh} prints the gate verdict and the \texttt{\#print\ axioms} report for all 36 required
declarations, and
\texttt{bash\ tests/\allowbreak tamper.\allowbreak sh} runs the eight-test tamper suite; the two \texttt{.\allowbreak bak-r26} files are the pre-promotion
copies, kept so the round-26 edits of \S{}4 can be audited. The paper-side material --- the layer reports,
the review rounds and their triages, the layer-13/14/15 comparisons, the mutation, contradiction and
satisfiability reports, and \texttt{namegate.\allowbreak py} --- is published in the \texttt{PaperL} folder of
\href{https://github.com/jpbald93/LEAN-verification}{github.com/jpbald93/LEAN-verification}, arranged as
\texttt{preprint/\allowbreak }, \texttt{reports/\allowbreak }, \texttt{reviews/\allowbreak }, \texttt{layer0/\allowbreak }, \texttt{layer13/\allowbreak }, \texttt{check/\allowbreak } and \texttt{evidence/\allowbreak }; Appendix A gives
the mapping from the artifact-relative paths this paper uses.

\hypertarget{the-36-gate-declarations}{%
\subsection{The 36 gate declarations}\label{the-36-gate-declarations}}

The gate (\texttt{gate.\allowbreak sh}) lists 36 required declarations; the anchor table (\texttt{layer11/\allowbreak ANCHORS.\allowbreak md}) resolves
all 36 to file:line, with kind and axiom list, by parsing the gate list, \texttt{Clifford/\allowbreak BridgeL10.\allowbreak lean} and
the two layer reports. Every one depends on \textbf{std} only.

{\scriptsize\begin{longtable}[]{@{}
  >{\raggedright\arraybackslash}p{(\columnwidth - 6\tabcolsep) * \real{0.2500}}
  >{\raggedright\arraybackslash}p{(\columnwidth - 6\tabcolsep) * \real{0.2500}}
  >{\raggedright\arraybackslash}p{(\columnwidth - 6\tabcolsep) * \real{0.2500}}
  >{\raggedright\arraybackslash}p{(\columnwidth - 6\tabcolsep) * \real{0.2500}}@{}}
\toprule\noalign{}
\begin{minipage}[b]{\linewidth}\raggedright
\#
\end{minipage} & \begin{minipage}[b]{\linewidth}\raggedright
declaration
\end{minipage} & \begin{minipage}[b]{\linewidth}\raggedright
file:line
\end{minipage} & \begin{minipage}[b]{\linewidth}\raggedright
kind
\end{minipage} \\
\midrule\noalign{}
\endhead
\bottomrule\noalign{}
\endlastfoot
1 & \texttt{Clifford.\allowbreak clifford\_\allowbreak one} & \href{https://github.com/jpbald93/paperl_clifford_chain/blob/d15464f0f4f4c8848e7d3e63dd68a03ce45433ae/Clifford/Basic.lean\#L601}{\texttt{Clifford/\allowbreak Basic.lean}:601} & theorem \\
2 & \texttt{Clifford.\allowbreak clifford\_\allowbreak two} & \href{https://github.com/jpbald93/paperl_clifford_chain/blob/d15464f0f4f4c8848e7d3e63dd68a03ce45433ae/Clifford/Basic.lean\#L693}{\texttt{Clifford/\allowbreak Basic.lean}:693} & theorem \\
3 & \texttt{Clifford.\allowbreak second\_\allowbreak onCircleM\_\allowbreak left} & \href{https://github.com/jpbald93/paperl_clifford_chain/blob/d15464f0f4f4c8848e7d3e63dd68a03ce45433ae/Clifford/Basic.lean\#L445}{\texttt{Clifford/\allowbreak Basic.lean}:445} & theorem \\
4 & \texttt{Clifford.\allowbreak second\_\allowbreak onCircleM\_\allowbreak right} & \href{https://github.com/jpbald93/paperl_clifford_chain/blob/d15464f0f4f4c8848e7d3e63dd68a03ce45433ae/Clifford/Basic.lean\#L449}{\texttt{Clifford/\allowbreak Basic.lean}:449} & theorem \\
5 & \texttt{Clifford.\allowbreak second\_\allowbreak eq\_\allowbreak coords} & \href{https://github.com/jpbald93/paperl_clifford_chain/blob/d15464f0f4f4c8848e7d3e63dd68a03ce45433ae/Clifford/Basic.lean\#L424}{\texttt{Clifford/\allowbreak Basic.lean}:424} & theorem \\
6 & \texttt{Clifford.\allowbreak onCircleM\_\allowbreak iff} & \href{https://github.com/jpbald93/paperl_clifford_chain/blob/d15464f0f4f4c8848e7d3e63dd68a03ce45433ae/Clifford/Basic.lean\#L385}{\texttt{Clifford/\allowbreak Basic.lean}:385} & theorem \\
7 & \texttt{Clifford.\allowbreak mD\_\allowbreak ne\_\allowbreak zero} & \href{https://github.com/jpbald93/paperl_clifford_chain/blob/d15464f0f4f4c8848e7d3e63dd68a03ce45433ae/Clifford/Basic.lean\#L520}{\texttt{Clifford/\allowbreak Basic.lean}:520} & theorem \\
8 & \texttt{Clifford.\allowbreak clifford\_\allowbreak two\_\allowbreak clean} & \href{https://github.com/jpbald93/paperl_clifford_chain/blob/d15464f0f4f4c8848e7d3e63dd68a03ce45433ae/Clifford/Clean.lean\#L123}{\texttt{Clifford/\allowbreak Clean.lean}:123} & theorem \\
9 & \texttt{Clifford.\allowbreak collinear\_\allowbreak of\_\allowbreak line} & \href{https://github.com/jpbald93/paperl_clifford_chain/blob/d15464f0f4f4c8848e7d3e63dd68a03ce45433ae/Clifford/Clean.lean\#L36}{\texttt{Clifford/\allowbreak Clean.lean}:36} & theorem \\
10 & \texttt{Clifford.\allowbreak miquel\_\allowbreak all\_\allowbreak eq} & \href{https://github.com/jpbald93/paperl_clifford_chain/blob/d15464f0f4f4c8848e7d3e63dd68a03ce45433ae/Clifford/Clean.lean\#L93}{\texttt{Clifford/\allowbreak Clean.lean}:93} & theorem \\
11 & \texttt{Clifford.\allowbreak clifford\_\allowbreak points\_\allowbreak on} & \href{https://github.com/jpbald93/paperl_clifford_chain/blob/d15464f0f4f4c8848e7d3e63dd68a03ce45433ae/Clifford/Clean.lean\#L43}{\texttt{Clifford/\allowbreak Clean.lean}:43} & theorem \\
12 & \texttt{Clifford.\allowbreak mem\_\allowbreak line\_\allowbreak of} & \href{https://github.com/jpbald93/paperl_clifford_chain/blob/d15464f0f4f4c8848e7d3e63dd68a03ce45433ae/Clifford/Clean.lean\#L24}{\texttt{Clifford/\allowbreak Clean.lean}:24} & theorem \\
13 & \texttt{Clifford.\allowbreak cliffordChain\_\allowbreak four} & \href{https://github.com/jpbald93/paperl_clifford_chain/blob/d15464f0f4f4c8848e7d3e63dd68a03ce45433ae/Clifford/Clean.lean\#L226}{\texttt{Clifford/\allowbreak Clean.lean}:226} & theorem \\
14 & \texttt{Clifford.\allowbreak cliffordChain\_\allowbreak five} & \href{https://github.com/jpbald93/paperl_clifford_chain/blob/d15464f0f4f4c8848e7d3e63dd68a03ce45433ae/Clifford/Clean.lean\#L245}{\texttt{Clifford/\allowbreak Clean.lean}:245} & theorem \\
15 & \texttt{Clifford.\allowbreak collinear\_\allowbreak iff\_\allowbreak conj\_\allowbreak mul} & \href{https://github.com/jpbald93/paperl_clifford_chain/blob/d15464f0f4f4c8848e7d3e63dd68a03ce45433ae/Clifford/Bridge.lean\#L22}{\texttt{Clifford/\allowbreak Bridge.lean}:22} & theorem \\
16 & \texttt{Clifford.\allowbreak transversal\_\allowbreak iff} & \href{https://github.com/jpbald93/paperl_clifford_chain/blob/d15464f0f4f4c8848e7d3e63dd68a03ce45433ae/Clifford/Bridge.lean\#L58}{\texttt{Clifford/\allowbreak Bridge.lean}:58} & theorem \\
17 & \texttt{Clifford.\allowbreak triP\_\allowbreak ne\_\allowbreak zero\_\allowbreak iff} & \href{https://github.com/jpbald93/paperl_clifford_chain/blob/d15464f0f4f4c8848e7d3e63dd68a03ce45433ae/Clifford/Bridge.lean\#L62}{\texttt{Clifford/\allowbreak Bridge.lean}:62} & theorem \\
18 & \texttt{Clifford.\allowbreak cross\_\allowbreak second} & \href{https://github.com/jpbald93/paperl_clifford_chain/blob/d15464f0f4f4c8848e7d3e63dd68a03ce45433ae/Clifford/Bridge.lean\#L70}{\texttt{Clifford/\allowbreak Bridge.lean}:70} & theorem \\
19 & \texttt{Clifford.\allowbreak cycP\_\allowbreak ne\_\allowbreak zero\_\allowbreak iff} & \href{https://github.com/jpbald93/paperl_clifford_chain/blob/d15464f0f4f4c8848e7d3e63dd68a03ce45433ae/Clifford/Bridge.lean\#L83}{\texttt{Clifford/\allowbreak Bridge.lean}:83} & theorem \\
20 & \texttt{Clifford.\allowbreak onCircleM\_\allowbreak iff\_\allowbreak mem\_\allowbreak sphere} & \href{https://github.com/jpbald93/paperl_clifford_chain/blob/d15464f0f4f4c8848e7d3e63dd68a03ce45433ae/Clifford/Bridge.lean\#L104}{\texttt{Clifford/\allowbreak Bridge.lean}:104} & theorem \\
21 & \texttt{Clifford.\allowbreak genTriple\_\allowbreak iff} & \href{https://github.com/jpbald93/paperl_clifford_chain/blob/d15464f0f4f4c8848e7d3e63dd68a03ce45433ae/Clifford/Bridge.lean\#L110}{\texttt{Clifford/\allowbreak Bridge.lean}:110} & theorem \\
22 & \texttt{Clifford.\allowbreak cD\_\allowbreak eq\_\allowbreak zero} & \href{https://github.com/jpbald93/paperl_clifford_chain/blob/d15464f0f4f4c8848e7d3e63dd68a03ce45433ae/Clifford/General.lean\#L32}{\texttt{Clifford/\allowbreak General.lean}:32} & theorem \\
23 & \texttt{Clifford.\allowbreak chain4\_\allowbreak cD} & \href{https://github.com/jpbald93/paperl_clifford_chain/blob/d15464f0f4f4c8848e7d3e63dd68a03ce45433ae/Clifford/General.lean\#L61}{\texttt{Clifford/\allowbreak General.lean}:61} & theorem \\
24 & \texttt{Clifford.\allowbreak chain5\_\allowbreak mem} & \href{https://github.com/jpbald93/paperl_clifford_chain/blob/d15464f0f4f4c8848e7d3e63dd68a03ce45433ae/Clifford/General.lean\#L87}{\texttt{Clifford/\allowbreak General.lean}:87} & theorem \\
25 & \texttt{Clifford.\allowbreak Counter6.\allowbreak cx\_\allowbreak genPosN} & \href{https://github.com/jpbald93/paperl_clifford_chain/blob/d15464f0f4f4c8848e7d3e63dd68a03ce45433ae/Clifford/Counter6.lean\#L33}{\texttt{Clifford/\allowbreak Counter6.lean}:33} & theorem \\
26 & \texttt{Clifford.\allowbreak Counter6.\allowbreak not\_\allowbreak chain6} & \href{https://github.com/jpbald93/paperl_clifford_chain/blob/d15464f0f4f4c8848e7d3e63dd68a03ce45433ae/Clifford/Counter6.lean\#L182}{\texttt{Clifford/\allowbreak Counter6.lean}:182} & theorem \\
27 & \texttt{Clifford.\allowbreak Counter6.\allowbreak not\_\allowbreak cliffordChain\_\allowbreak six} & \href{https://github.com/jpbald93/paperl_clifford_chain/blob/d15464f0f4f4c8848e7d3e63dd68a03ce45433ae/Clifford/Counter6.lean\#L227}{\texttt{Clifford/\allowbreak Counter6.lean}:227} & theorem \\
28 & \texttt{Clifford.\allowbreak cliffordChain\_\allowbreak three} & \href{https://github.com/jpbald93/paperl_clifford_chain/blob/d15464f0f4f4c8848e7d3e63dd68a03ce45433ae/Clifford/Classify.lean\#L43}{\texttt{Clifford/\allowbreak Classify.lean}:43} & theorem \\
29 & \texttt{Clifford.\allowbreak not\_\allowbreak chain\_\allowbreak of\_\allowbreak extends} & \href{https://github.com/jpbald93/paperl_clifford_chain/blob/d15464f0f4f4c8848e7d3e63dd68a03ce45433ae/Clifford/Classify.lean\#L57}{\texttt{Clifford/\allowbreak Classify.lean}:57} & theorem \\
30 & \texttt{Clifford.\allowbreak not\_\allowbreak cliffordChain\_\allowbreak of\_\allowbreak six\_\allowbreak le} & \href{https://github.com/jpbald93/paperl_clifford_chain/blob/d15464f0f4f4c8848e7d3e63dd68a03ce45433ae/Clifford/Classify.lean\#L219}{\texttt{Clifford/\allowbreak Classify.lean}:219} & theorem \\
31 & \texttt{Clifford.\allowbreak cliffordChain\_\allowbreak iff} & \href{https://github.com/jpbald93/paperl_clifford_chain/blob/d15464f0f4f4c8848e7d3e63dd68a03ce45433ae/Clifford/Classify.lean\#L227}{\texttt{Clifford/\allowbreak Classify.lean}:227} & theorem \\
32 & \texttt{Clifford.\allowbreak MorleyKernel.\allowbreak cliffordChainS\_\allowbreak inf} & \href{https://github.com/jpbald93/paperl_clifford_chain/blob/d15464f0f4f4c8848e7d3e63dd68a03ce45433ae/Clifford/MorleyKernel.lean\#L5592}{\texttt{Clifford/\allowbreak MorleyKernel.lean}:5592} & theorem \\
33 & \texttt{Clifford.\allowbreak MorleyKernel.\allowbreak heven\_\allowbreak inf\_\allowbreak all} & \href{https://github.com/jpbald93/paperl_clifford_chain/blob/d15464f0f4f4c8848e7d3e63dd68a03ce45433ae/Clifford/MorleyKernel.lean\#L5586}{\texttt{Clifford/\allowbreak MorleyKernel.lean}:5586} & theorem \\
34 & \texttt{Clifford.\allowbreak MorleyKernel.\allowbreak EvenCorner\_\allowbreak proof} & \href{https://github.com/jpbald93/paperl_clifford_chain/blob/d15464f0f4f4c8848e7d3e63dd68a03ce45433ae/Clifford/MorleyKernel.lean\#L5584}{\texttt{Clifford/\allowbreak MorleyKernel.lean}:5584} & theorem \\
35 & \texttt{Clifford.\allowbreak MorleyKernel.\allowbreak evenCornerQL\_\allowbreak proof} & \href{https://github.com/jpbald93/paperl_clifford_chain/blob/d15464f0f4f4c8848e7d3e63dd68a03ce45433ae/Clifford/MorleyKernel.lean\#L5539}{\texttt{Clifford/\allowbreak MorleyKernel.lean}:5539} & theorem \\
36 & \texttt{Clifford.\allowbreak GenPosN.\allowbreak exists\_\allowbreak snoc} & \href{https://github.com/jpbald93/paperl_clifford_chain/blob/d15464f0f4f4c8848e7d3e63dd68a03ce45433ae/Clifford/Classify.lean\#L167}{\texttt{Clifford/\allowbreak Classify.lean}:167} & theorem \\
\end{longtable}}

(\texttt{GenPosN.\allowbreak exists\_\allowbreak snoc} is namespace-qualified in its declaration, which is why a naive
\texttt{theorem\ \textless{}short\textgreater{}\ match} misses it; the anchor harvester records that explicitly.)

\hypertarget{the-19-layer-10-bridges}{%
\subsection{The 19 layer-10 bridges}\label{the-19-layer-10-bridges}}

\texttt{Clifford/\allowbreak BridgeL10.\allowbreak lean} adds 19 theorems, each with a docstring, all in \texttt{namespace\ Clifford}, all
on \textbf{std}, none of them in the gate's REQUIRED list. They tie the paper's own notions to Mathlib:

{\scriptsize\begin{longtable}[]{@{}
  >{\raggedright\arraybackslash}p{(\columnwidth - 6\tabcolsep) * \real{0.2500}}
  >{\raggedright\arraybackslash}p{(\columnwidth - 6\tabcolsep) * \real{0.2500}}
  >{\raggedright\arraybackslash}p{(\columnwidth - 6\tabcolsep) * \real{0.2500}}
  >{\raggedright\arraybackslash}p{(\columnwidth - 6\tabcolsep) * \real{0.2500}}@{}}
\toprule\noalign{}
\begin{minipage}[b]{\linewidth}\raggedright
\#
\end{minipage} & \begin{minipage}[b]{\linewidth}\raggedright
declaration
\end{minipage} & \begin{minipage}[b]{\linewidth}\raggedright
file:line
\end{minipage} & \begin{minipage}[b]{\linewidth}\raggedright
notion $\to$ Mathlib counterpart
\end{minipage} \\
\midrule\noalign{}
\endhead
\bottomrule\noalign{}
\endlastfoot
L10.1 & \texttt{det3\_\allowbreak eq\_\allowbreak det} & \href{https://github.com/jpbald93/paperl_clifford_chain/blob/d15464f0f4f4c8848e7d3e63dd68a03ce45433ae/Clifford/BridgeL10.lean\#L32}{\texttt{Clifford/\allowbreak BridgeL10.lean}:32} & \texttt{det3} $\to$ \texttt{Matrix.\allowbreak det} (Fin 3) \\
L10.2 & \texttt{crs\_\allowbreak eq\_\allowbreak det2} & \href{https://github.com/jpbald93/paperl_clifford_chain/blob/d15464f0f4f4c8848e7d3e63dd68a03ce45433ae/Clifford/BridgeL10.lean\#L38}{\texttt{Clifford/\allowbreak BridgeL10.lean}:38} & \texttt{crs} $\to$ \texttt{Matrix.\allowbreak det} (Fin 2) \\
L10.3 & \texttt{det\_\allowbreak fin\_\allowbreak four\_\allowbreak eq\_\allowbreak det4} & \href{https://github.com/jpbald93/paperl_clifford_chain/blob/d15464f0f4f4c8848e7d3e63dd68a03ce45433ae/Clifford/BridgeL10.lean\#L43}{\texttt{Clifford/\allowbreak BridgeL10.lean}:43} & any \texttt{Matrix\ (Fin\ 4)\ (Fin\ 4)\ R} $\to$ \texttt{det4} \\
L10.4 & \texttt{det\_\allowbreak fin\_\allowbreak five\_\allowbreak eq\_\allowbreak det5} & \href{https://github.com/jpbald93/paperl_clifford_chain/blob/d15464f0f4f4c8848e7d3e63dd68a03ce45433ae/Clifford/BridgeL10.lean\#L53}{\texttt{Clifford/\allowbreak BridgeL10.lean}:53} & any \texttt{Matrix\ (Fin\ 5)\ (Fin\ 5)\ R} $\to$ \texttt{det5} \\
L10.5 & \texttt{det4\_\allowbreak eq\_\allowbreak det} & \href{https://github.com/jpbald93/paperl_clifford_chain/blob/d15464f0f4f4c8848e7d3e63dd68a03ce45433ae/Clifford/BridgeL10.lean\#L63}{\texttt{Clifford/\allowbreak BridgeL10.lean}:63} & \texttt{det4} $\to$ \texttt{Matrix.\allowbreak det} \\
L10.6 & \texttt{det5\_\allowbreak eq\_\allowbreak det} & \href{https://github.com/jpbald93/paperl_clifford_chain/blob/d15464f0f4f4c8848e7d3e63dd68a03ce45433ae/Clifford/BridgeL10.lean\#L70}{\texttt{Clifford/\allowbreak BridgeL10.lean}:70} & \texttt{det5} $\to$ \texttt{Matrix.\allowbreak det} \\
L10.7 & \texttt{mN\_\allowbreak eq\_\allowbreak det} & \href{https://github.com/jpbald93/paperl_clifford_chain/blob/d15464f0f4f4c8848e7d3e63dd68a03ce45433ae/Clifford/BridgeL10.lean\#L78}{\texttt{Clifford/\allowbreak BridgeL10.lean}:78} & \texttt{mN} $\to$ \texttt{Matrix.\allowbreak det} \\
L10.8 & \texttt{mD\_\allowbreak eq\_\allowbreak det} & \href{https://github.com/jpbald93/paperl_clifford_chain/blob/d15464f0f4f4c8848e7d3e63dd68a03ce45433ae/Clifford/BridgeL10.lean\#L86}{\texttt{Clifford/\allowbreak BridgeL10.lean}:86} & \texttt{mD} $\to$ \texttt{Matrix.\allowbreak det} \\
L10.9 & \texttt{crossVec\_\allowbreak eq\_\allowbreak crossProduct} & \href{https://github.com/jpbald93/paperl_clifford_chain/blob/d15464f0f4f4c8848e7d3e63dd68a03ce45433ae/Clifford/BridgeL10.lean\#L95}{\texttt{Clifford/\allowbreak BridgeL10.lean}:95} & \texttt{crossVec} (2$\times$3 matrices only) $\to$ \texttt{Mathlib} \texttt{crossProduct} \\
L10.10 & \texttt{miquel\_\allowbreak eq\_\allowbreak det\_\allowbreak div} & \href{https://github.com/jpbald93/paperl_clifford_chain/blob/d15464f0f4f4c8848e7d3e63dd68a03ce45433ae/Clifford/BridgeL10.lean\#L105}{\texttt{Clifford/\allowbreak BridgeL10.lean}:105} & \texttt{miquel} $\to$ ratio of two \texttt{Matrix.\allowbreak det} \\
L10.11 & \texttt{triP\_\allowbreak eq\_\allowbreak zero\_\allowbreak iff\_\allowbreak collinear} & \href{https://github.com/jpbald93/paperl_clifford_chain/blob/d15464f0f4f4c8848e7d3e63dd68a03ce45433ae/Clifford/BridgeL10.lean\#L117}{\texttt{Clifford/\allowbreak BridgeL10.lean}:117} & \texttt{triP\ =\ 0} $\leftrightarrow$ \texttt{Collinear\ ℝ} \\
L10.12 & \texttt{colD\_\allowbreak eq\_\allowbreak zero\_\allowbreak iff\_\allowbreak collinear} & \href{https://github.com/jpbald93/paperl_clifford_chain/blob/d15464f0f4f4c8848e7d3e63dd68a03ce45433ae/Clifford/BridgeL10.lean\#L131}{\texttt{Clifford/\allowbreak BridgeL10.lean}:131} & \texttt{colD\ =\ 0} $\leftrightarrow$ \texttt{Collinear\ ℝ} \\
L10.13 & \texttt{cD\_\allowbreak eq\_\allowbreak det} & \href{https://github.com/jpbald93/paperl_clifford_chain/blob/d15464f0f4f4c8848e7d3e63dd68a03ce45433ae/Clifford/BridgeL10.lean\#L137}{\texttt{Clifford/\allowbreak BridgeL10.lean}:137} & \texttt{cD} $\to$ \texttt{Matrix.\allowbreak det} \\
L10.14 & \texttt{cospherical\_\allowbreak or\_\allowbreak collinear\_\allowbreak of\_\allowbreak eqn} & \href{https://github.com/jpbald93/paperl_clifford_chain/blob/d15464f0f4f4c8848e7d3e63dd68a03ce45433ae/Clifford/BridgeL10.lean\#L145}{\texttt{Clifford/\allowbreak BridgeL10.lean}:145} & from an explicit equation to \texttt{Cospherical\ ∨\ Collinear} \\
L10.15 & \texttt{exists\_\allowbreak eqn\_\allowbreak of\_\allowbreak cD} & \href{https://github.com/jpbald93/paperl_clifford_chain/blob/d15464f0f4f4c8848e7d3e63dd68a03ce45433ae/Clifford/BridgeL10.lean\#L189}{\texttt{Clifford/\allowbreak BridgeL10.lean}:189} & \texttt{cD} $\to$ real linear equation \texttt{A\textbar{}p\textbar{}$^2$+Bp.re+Cp.im+D=0} \\
L10.16 & \texttt{cD\_\allowbreak eq\_\allowbreak zero\_\allowbreak iff} & \href{https://github.com/jpbald93/paperl_clifford_chain/blob/d15464f0f4f4c8848e7d3e63dd68a03ce45433ae/Clifford/BridgeL10.lean\#L209}{\texttt{Clifford/\allowbreak BridgeL10.lean}:209} & \texttt{cD\ =\ 0} $\leftrightarrow$ \texttt{Cospherical\ ∨\ Collinear} (no extra hypothesis) \\
L10.17 & \texttt{cycP\_\allowbreak eq\_\allowbreak cD} & \href{https://github.com/jpbald93/paperl_clifford_chain/blob/d15464f0f4f4c8848e7d3e63dd68a03ce45433ae/Clifford/BridgeL10.lean\#L221}{\texttt{Clifford/\allowbreak BridgeL10.lean}:221} & \texttt{cycP\ =\ −2·I·cD\ 0\ a\ b\ c} \\
L10.18 & \texttt{cycP\_\allowbreak eq\_\allowbreak zero\_\allowbreak iff} & \href{https://github.com/jpbald93/paperl_clifford_chain/blob/d15464f0f4f4c8848e7d3e63dd68a03ce45433ae/Clifford/BridgeL10.lean\#L239}{\texttt{Clifford/\allowbreak BridgeL10.lean}:239} & \texttt{cycP\ =\ 0} $\leftrightarrow$ \texttt{Cospherical\ ∨\ Collinear} of \texttt{\{0,a,b,c\}} \\
L10.19 & \texttt{onCircle3\_\allowbreak iff\_\allowbreak cospherical} & \href{https://github.com/jpbald93/paperl_clifford_chain/blob/d15464f0f4f4c8848e7d3e63dd68a03ce45433ae/Clifford/BridgeL10.lean\#L248}{\texttt{Clifford/\allowbreak BridgeL10.lean}:248} & \texttt{OnCircle3} $\leftrightarrow$ \texttt{Cospherical} (definitional) \\
\end{longtable}}

Two points deserve to be in the paper rather than only in the report. First, \texttt{cD\_\allowbreak eq\_\allowbreak zero\_\allowbreak iff} needs
\textbf{no} nondegeneracy hypothesis: degenerate point sets fall under one of the two disjuncts, and the
\texttt{←} direction is \texttt{Clean.\allowbreak exists\_\allowbreak gcircle} plus \texttt{General.\allowbreak cD\_\allowbreak eq\_\allowbreak zero}, while \texttt{→} goes through a nonzero
kernel vector, a real equation, and a circle if \texttt{A\ ≠\ 0} and a line if \texttt{A\ =\ 0}. This is the lemma
that ties \texttt{cD\ =\ 0} to the standard \texttt{Cospherical\ ∨\ Collinear} reading of ``four points on a circle''; it
does not by itself bridge the sphere predicate \texttt{IsGCircleS} (\S{}5.6). Second, \texttt{cycP\_\allowbreak eq\_\allowbreak cD}: the first
sign guess \texttt{+2·I} was false (\texttt{ring} failed with goal \texttt{LHS\ =\ −LHS}) and the correct factor is \texttt{−2·I};
the report records that, and the draft should not hide it --- a formalisation account that reports only
its successes is not a verification record.

\hypertarget{the-statement-declarations-behind-2}{%
\subsection{The statement declarations behind \S{}2}\label{the-statement-declarations-behind-2}}

These are the declarations the prose of \S{}1--2 leans on and that the 36-row gate list does not cover;
each is checked at the cited line by the same layer-11 script that checks the prose citations.

{\scriptsize\begin{longtable}[]{@{}
  >{\raggedright\arraybackslash}p{(\columnwidth - 6\tabcolsep) * \real{0.2500}}
  >{\raggedright\arraybackslash}p{(\columnwidth - 6\tabcolsep) * \real{0.2500}}
  >{\raggedright\arraybackslash}p{(\columnwidth - 6\tabcolsep) * \real{0.2500}}
  >{\raggedright\arraybackslash}p{(\columnwidth - 6\tabcolsep) * \real{0.2500}}@{}}
\toprule\noalign{}
\begin{minipage}[b]{\linewidth}\raggedright
\#
\end{minipage} & \begin{minipage}[b]{\linewidth}\raggedright
declaration
\end{minipage} & \begin{minipage}[b]{\linewidth}\raggedright
file:line
\end{minipage} & \begin{minipage}[b]{\linewidth}\raggedright
role in \S{}1--2
\end{minipage} \\
\midrule\noalign{}
\endhead
\bottomrule\noalign{}
\endlastfoot
S1 & \texttt{second} & \href{https://github.com/jpbald93/paperl_clifford_chain/blob/d15464f0f4f4c8848e7d3e63dd68a03ce45433ae/Clifford/Basic.lean\#L394}{\texttt{Clifford/\allowbreak Basic.lean}:394} & second intersection of two circles \\
S2 & \texttt{Transversal} & \href{https://github.com/jpbald93/paperl_clifford_chain/blob/d15464f0f4f4c8848e7d3e63dd68a03ce45433ae/Clifford/Basic.lean\#L398}{\texttt{Clifford/\allowbreak Basic.lean}:398} & pairwise non-collinearity of \texttt{M,\ c,\ d} \\
S3 & \texttt{GenTriple} & \href{https://github.com/jpbald93/paperl_clifford_chain/blob/d15464f0f4f4c8848e7d3e63dd68a03ce45433ae/Clifford/Basic.lean\#L402}{\texttt{Clifford/\allowbreak Basic.lean}:402} & the two conjuncts of triple general position \\
S4 & \texttt{lineS} & \href{https://github.com/jpbald93/paperl_clifford_chain/blob/d15464f0f4f4c8848e7d3e63dd68a03ce45433ae/Clifford/Mobius.lean\#L24}{\texttt{Clifford/\allowbreak Mobius.lean}:24} & a line plus \texttt{∞} \\
S5 & \texttt{IsGCircleS} & \href{https://github.com/jpbald93/paperl_clifford_chain/blob/d15464f0f4f4c8848e7d3e63dd68a03ce45433ae/Clifford/Mobius.lean\#L28}{\texttt{Clifford/\allowbreak Mobius.lean}:28} & generalized circles on the sphere \\
S6 & \texttt{ChainS} & \href{https://github.com/jpbald93/paperl_clifford_chain/blob/d15464f0f4f4c8848e7d3e63dd68a03ce45433ae/Clifford/Mobius.lean\#L35}{\texttt{Clifford/\allowbreak Mobius.lean}:35} & the sphere recursion \\
S7 & \texttt{CliffordChainS} & \href{https://github.com/jpbald93/paperl_clifford_chain/blob/d15464f0f4f4c8848e7d3e63dd68a03ce45433ae/Clifford/Mobius.lean\#L46}{\texttt{Clifford/\allowbreak Mobius.lean}:46} & existence of a level-n object \\
S8 & \texttt{chainS\_\allowbreak of\_\allowbreak chain} & \href{https://github.com/jpbald93/paperl_clifford_chain/blob/d15464f0f4f4c8848e7d3e63dd68a03ce45433ae/Clifford/Mobius.lean\#L119}{\texttt{Clifford/\allowbreak Mobius.lean}:119} & transport of a finite chain object \\
S9 & \texttt{chainS\_\allowbreak of\_\allowbreak chain\_\allowbreak pt} & \href{https://github.com/jpbald93/paperl_clifford_chain/blob/d15464f0f4f4c8848e7d3e63dd68a03ce45433ae/Clifford/Mobius.lean\#L152}{\texttt{Clifford/\allowbreak Mobius.lean}:152} & transport of a singleton chain object \\
S10 & \texttt{cliffordChainS\_\allowbreak of\_\allowbreak le\_\allowbreak five} & \href{https://github.com/jpbald93/paperl_clifford_chain/blob/d15464f0f4f4c8848e7d3e63dd68a03ce45433ae/Clifford/Mobius.lean\#L178}{\texttt{Clifford/\allowbreak Mobius.lean}:178} & the sphere base cases \texttt{n\ ≤\ 5} \\
S11 & \texttt{cliffordChainS\_\allowbreak all\_\allowbreak of} & \href{https://github.com/jpbald93/paperl_clifford_chain/blob/d15464f0f4f4c8848e7d3e63dd68a03ce45433ae/Clifford/Mobius.lean\#L330}{\texttt{Clifford/\allowbreak Mobius.lean}:330} & the two-step induction to \texttt{ChainS} \\
S12 & \texttt{chainS\_\allowbreak six} & \href{https://github.com/jpbald93/paperl_clifford_chain/blob/d15464f0f4f4c8848e7d3e63dd68a03ce45433ae/Clifford/Mobius.lean\#L651}{\texttt{Clifford/\allowbreak Mobius.lean}:651} & level-6 repair of the six-centre witness \\
S13 & \texttt{chainS\_\allowbreak six\_\allowbreak unique} & \href{https://github.com/jpbald93/paperl_clifford_chain/blob/d15464f0f4f4c8848e7d3e63dd68a03ce45433ae/Clifford/Mobius.lean\#L734}{\texttt{Clifford/\allowbreak Mobius.lean}:734} & uniqueness there (this configuration only) \\
S14 & \texttt{chainS\_\allowbreak six\_\allowbreak iff} & \href{https://github.com/jpbald93/paperl_clifford_chain/blob/d15464f0f4f4c8848e7d3e63dd68a03ce45433ae/Clifford/Mobius.lean\#L779}{\texttt{Clifford/\allowbreak Mobius.lean}:779} & the iff form of S12--S13 \\
S15 & \texttt{crossVec} & \href{https://github.com/jpbald93/paperl_clifford_chain/blob/d15464f0f4f4c8848e7d3e63dd68a03ce45433ae/Clifford/MorleyKernel.lean\#L159}{\texttt{Clifford/\allowbreak MorleyKernel.lean}:159} & signed maximal minors \\
S16 & \texttt{KP} & \href{https://github.com/jpbald93/paperl_clifford_chain/blob/d15464f0f4f4c8848e7d3e63dd68a03ce45433ae/Clifford/MorleyKernel.lean\#L214}{\texttt{Clifford/\allowbreak MorleyKernel.lean}:214} & row-scaled Morley matrix \\
S17 & \texttt{levelPointN} & \href{https://github.com/jpbald93/paperl_clifford_chain/blob/d15464f0f4f4c8848e7d3e63dd68a03ce45433ae/Clifford/MorleyKernel.lean\#L452}{\texttt{Clifford/\allowbreak MorleyKernel.lean}:452} & \texttt{$-$v\_last\ /\ v\_\{n/2\}} \\
S18 & \texttt{EptInf} & \href{https://github.com/jpbald93/paperl_clifford_chain/blob/d15464f0f4f4c8848e7d3e63dd68a03ce45433ae/Clifford/MorleyKernel.lean\#L2549}{\texttt{Clifford/\allowbreak MorleyKernel.lean}:2549} & the totalized projective level point \\
S19 & \texttt{h4\_\allowbreak pt\_\allowbreak inf} & \href{https://github.com/jpbald93/paperl_clifford_chain/blob/d15464f0f4f4c8848e7d3e63dd68a03ce45433ae/Clifford/MorleyKernel.lean\#L2566}{\texttt{Clifford/\allowbreak MorleyKernel.lean}:2566} & the level-4 certificate for \texttt{EptInf} \\
S20 & \texttt{parent\_\allowbreak crossVec\_\allowbreak zero\_\allowbreak of\_\allowbreak mid} & \href{https://github.com/jpbald93/paperl_clifford_chain/blob/d15464f0f4f4c8848e7d3e63dd68a03ce45433ae/Clifford/MorleyKernel.lean\#L4998}{\texttt{Clifford/\allowbreak MorleyKernel.lean}:4998} & \texttt{V\ =\ 0} in the corner \\
S21 & \texttt{EvenCorner} & \href{https://github.com/jpbald93/paperl_clifford_chain/blob/d15464f0f4f4c8848e7d3e63dd68a03ce45433ae/Clifford/MorleyKernel.lean\#L5253}{\texttt{Clifford/\allowbreak MorleyKernel.lean}:5253} & the geometric corner statement \\
S22 & \texttt{evenCorner\_\allowbreak of\_\allowbreak QL} & \href{https://github.com/jpbald93/paperl_clifford_chain/blob/d15464f0f4f4c8848e7d3e63dd68a03ce45433ae/Clifford/MorleyKernel.lean\#L5287}{\texttt{Clifford/\allowbreak MorleyKernel.lean}:5287} & corner from the scalar identity \\
S23 & \texttt{EvenCornerQL} & \href{https://github.com/jpbald93/paperl_clifford_chain/blob/d15464f0f4f4c8848e7d3e63dd68a03ce45433ae/Clifford/MorleyKernel.lean\#L5317}{\texttt{Clifford/\allowbreak MorleyKernel.lean}:5317} & the scalar identity \\
S24 & \texttt{evenCorner\_\allowbreak of\_\allowbreak cornerQL} & \href{https://github.com/jpbald93/paperl_clifford_chain/blob/d15464f0f4f4c8848e7d3e63dd68a03ce45433ae/Clifford/MorleyKernel.lean\#L5325}{\texttt{Clifford/\allowbreak MorleyKernel.lean}:5325} & packaging \\
S25 & \texttt{evenCornerQL\_\allowbreak proof} & \href{https://github.com/jpbald93/paperl_clifford_chain/blob/d15464f0f4f4c8848e7d3e63dd68a03ce45433ae/Clifford/MorleyKernel.lean\#L5539}{\texttt{Clifford/\allowbreak MorleyKernel.lean}:5539} & round-26 proof of S23 \\
S26 & \texttt{EvenCorner\_\allowbreak proof} & \href{https://github.com/jpbald93/paperl_clifford_chain/blob/d15464f0f4f4c8848e7d3e63dd68a03ce45433ae/Clifford/MorleyKernel.lean\#L5584}{\texttt{Clifford/\allowbreak MorleyKernel.lean}:5584} & S21 from S25 \\
S27 & \texttt{heven\_\allowbreak inf\_\allowbreak all} & \href{https://github.com/jpbald93/paperl_clifford_chain/blob/d15464f0f4f4c8848e7d3e63dd68a03ce45433ae/Clifford/MorleyKernel.lean\#L5586}{\texttt{Clifford/\allowbreak MorleyKernel.lean}:5586} & the even payload \\
S28 & \texttt{hodd\_inf\_all\textquotesingle{}} & \href{https://github.com/jpbald93/paperl_clifford_chain/blob/d15464f0f4f4c8848e7d3e63dd68a03ce45433ae/Clifford/MorleyKernel.lean\#L4643}{\texttt{Clifford/\allowbreak MorleyKernel.lean}:4643} & the odd payload \\
S29 & \texttt{cliffordChainS\_\allowbreak inf} & \href{https://github.com/jpbald93/paperl_clifford_chain/blob/d15464f0f4f4c8848e7d3e63dd68a03ce45433ae/Clifford/MorleyKernel.lean\#L5592}{\texttt{Clifford/\allowbreak MorleyKernel.lean}:5592} & the headline, \texttt{n\ ≥\ 3} \\
S30 & \texttt{not\_\allowbreak cliffordChain\_\allowbreak of\_\allowbreak six\_\allowbreak le} & \href{https://github.com/jpbald93/paperl_clifford_chain/blob/d15464f0f4f4c8848e7d3e63dd68a03ce45433ae/Clifford/Classify.lean\#L219}{\texttt{Clifford/\allowbreak Classify.lean}:219} & finite-plane failure for \texttt{n\ ≥\ 6} \\
\end{longtable}}

The gate's earlier omission of \texttt{ChainS}, \texttt{CliffordChainS}, \texttt{lineS}, \texttt{IsGCircleS}, \texttt{chainS\_\allowbreak six*},
\texttt{evenCorner\_\allowbreak of\_\allowbreak QL} and \texttt{evenCorner\_\allowbreak of\_\allowbreak cornerQL} is what S1--S30 close; it is still a name-and-line
check, not a signature audit, and \texttt{L12\_\allowbreak REVIEW.\allowbreak md} did not perform one (it could not inspect those
types at the time). v0.3 quotes the definitions themselves in \S{}2.1, so a reader can compare the
paper's prose against the code without running Lean.


\hypertarget{the-verification-record}{%
\section{The verification record}\label{the-verification-record}}

All figures below are quoted from the recorded reports; the runner is the \texttt{leancheck} battery that
was developed for this paper series. Host: \texttt{jack}. Report
\texttt{\textasciitilde/\allowbreak leancheck/\allowbreak reports/\allowbreak paperL\_\allowbreak jack\_\allowbreak 20261119\_\allowbreak 020713.\allowbreak md} (+ \texttt{.\allowbreak log}); the verdict table is computed from
\texttt{\textasciitilde/\allowbreak leancheck/\allowbreak out/\allowbreak paperL\_\allowbreak jack/\allowbreak results.\allowbreak jsonl} (935 rows) and \texttt{\textasciitilde/\allowbreak lc\_\allowbreak gap\_\allowbreak paperL\_\allowbreak jack/\allowbreak results.\allowbreak jsonl}
(771 rows), independently recounted. The battery ran with \texttt{-\/\allowbreak -fast} against the promoted library at
\texttt{\textasciitilde/\allowbreak paperL\_\allowbreak jack}: \texttt{Clifford/\allowbreak MorleyKernel.\allowbreak lean} sha \texttt{dce7b011…}, gate REQUIRED 36. Wall clock:
02:07:13Z $\to$ 19:28:24Z on 2026-11-19, \textbf{17 h 21 min}, and the verdict line is
\texttt{LEANCHECK\ DONE\ Thu\ Nov\ 19\ 07\allowbreak :28\allowbreak :24\ PM\ UTC\ 2026}.

{\scriptsize\begin{longtable}[]{@{}
  >{\raggedright\arraybackslash}p{(\columnwidth - 4\tabcolsep) * \real{0.3333}}
  >{\raggedright\arraybackslash}p{(\columnwidth - 4\tabcolsep) * \real{0.3333}}
  >{\raggedright\arraybackslash}p{(\columnwidth - 4\tabcolsep) * \real{0.3333}}@{}}
\toprule\noalign{}
\begin{minipage}[b]{\linewidth}\raggedright
layer
\end{minipage} & \begin{minipage}[b]{\linewidth}\raggedright
what it checks
\end{minipage} & \begin{minipage}[b]{\linewidth}\raggedright
result
\end{minipage} \\
\midrule\noalign{}
\endhead
\bottomrule\noalign{}
\endlastfoot
1 & gate: every REQUIRED declaration present, \texttt{\#print\ axioms} $\subseteq$ std & \textbf{PASS (36 declarations, standard axioms only)} \\
2 & tampering: 8 hand-made cheat attempts must be rejected & \textbf{8/8 rejected} --- the log records \texttt{ALL\ TAMPERING\ TESTS\ REJECTED} \\
baseline & the 8 audited source files compile as given & \textbf{8/8 ok} \\
4 & mutation: inject a wrong constant at a declaration and require a downstream error & battery (2026-11-19, re-run 2026-11-25): \textbf{535 decisive kills + 95 timeout/heartbeat}, \textbf{1 SURVIVED}, \textbf{0 error-unattributed} over 259 declarations of 8 source files, only 61 of \texttt{MorleyKernel.\allowbreak lean}'s 343; late block (2026-11-25): \textbf{872 kills + 3 timeout}, \textbf{16 SURVIVED}, 0 unattributed over 257 further \texttt{MorleyKernel} declarations --- together \textbf{318 of the 343} in that file (see \S{}5.1) \\
5 & contradiction probes: try to prove \texttt{False} from a theorem's hypotheses with limited tactics & \textbf{0 CONTRADICTORY} over 551 of the 577 declarations (310 \texttt{ok}, 218 no-hypothesis, 22 timeout, 1 probe-invalid, 0 error-unattributed): the battery's 294 plus a late block over 257 further \texttt{MorleyKernel} declarations at lines 816--5586 (163 \texttt{ok}, 93 no-hypothesis, 1 probe-invalid). \texttt{ok} means the limited search failed, not that satisfiability was certified (\S{}5.2) \\
6 & ablation: delete a hypothesis and check whether the proof still compiles & \textbf{501 needed / 8 unused}, 0 skipped, 0 timeout, 0 bad-control, over 262 declarations \\
7 & linter (library files) & \textbf{0 warnings in 8 files}; the scratch leg was not run by the battery and is quoted from \texttt{layer10/\allowbreak REPORT\_\allowbreak L10.\allowbreak md} (\S{}4 below) \\
8 & satisfiability certification & \textbf{run on two machines, 2026-11-25/26}: jack \textbf{400 of 596 (67.1\%)}, john \textbf{394 of 584 (67.5\%)} certified; across the 577 declarations both ran, \textbf{397 (68.8\%)} are certified by at least one and 180 by neither; 0 stmt-invalid, nothing refuted (\S{}5.2) \\
8b & theorem $\leftrightarrow$ witness linkage & \textbf{0 links} --- neither tree contains hand-written \texttt{*\_\allowbreak sat} witnesses, so there is nothing to link (\S{}5.2) \\
9 & notion inventory: which notions have a Mathlib bridge & \textbf{137 notions, 24 flagged \texttt{has-bridge}, 113 not} --- a heuristic flag, not a semantic count (\S{}5.6) \\
10 & definition bridges (\texttt{BridgeL10.\allowbreak lean}) & 19 theorems; \textbf{13 of the 113} moved to bridged; 25 reclassified INTERNAL; \textbf{75 remain unbridged} (arithmetic, same caveat) \\
11 & name/line anchors for the draft & \textbf{PASS} --- 143 occurrences / 98 distinct citations and all 85 table rows resolve (\S{}3) \\
\end{longtable}}

The eight tamper tests are \texttt{sorry}, \texttt{ax2line}, \texttt{evalfake}, \texttt{macro}, \texttt{wrap} (nonstandard-axiom catch),
\texttt{falsehyp}, \texttt{cleanand}, \texttt{classify6}; the log reports \texttt{ALL\ TAMPERING\ TESTS\ REJECTED} (there is no
\texttt{TAMPER\_\allowbreak RC} string in it, and none is claimed here). The gate is not a build-time check alone: it
also pins the axiom set, which is why the two later edits in round 26b --- inlining two tactic macros
at all \textbf{13 use sites} (8 \texttt{fin\_\allowbreak idx} + 5 \texttt{fin\_\allowbreak idxO}; the 15 identifier occurrences include their two
definition lines), and rewording three comments containing the word ``prefix'' (module comments in the
promoted file, plus one docstring moved into a module comment) --- were forced by the gate rather than
by the mathematics (a \texttt{macro} could redefine \texttt{\#print\ axioms}; \texttt{\textbackslash{}bprefix\textbackslash{}b} is a forbidden token).
Both edits are recorded, and the 267-line round-26 proof is untouched by them.

The single mutation survivor of the 2026-11-19 battery is \texttt{cliffordChain\_\allowbreak one}; the late block's
sixteen are analysed in \S{}5.1. The injected mutant is exactly
\texttt{CliffordChain\ 1\ →\ CliffordChain\ 2} (row \texttt{lit1-\textgreater{}2@15}): a different but equally true statement, with
a proof term of the same shape, so it survives compilation of its own declaration and the
discrepancy surfaces downstream inside \texttt{cliffordChain\_\allowbreak iff} (\href{https://github.com/jpbald93/paperl_clifford_chain/blob/d15464f0f4f4c8848e7d3e63dd68a03ce45433ae/Clifford/Classify.lean\#L227}{\texttt{Clifford/\allowbreak Classify.lean}:227}), which needs the
\texttt{n\ =\ 1} instance and finds the wrong type. That is a limitation of mutation testing at a declaration
whose statement is level-parametric, not a hole in the proof; nothing in the battery classifies it as
either.

Build receipts at the pin. On jack: \texttt{lake\ build\ Clifford} rc = 0 (8,716 jobs at the layer-10
collection), \texttt{gate.\allowbreak sh} PASS (36 declarations, standard axioms only), \texttt{bash\ tests/\allowbreak tamper.\allowbreak sh} 8/8 rejected.
The scratch file \texttt{scratch/\allowbreak ChainEvenCornerQL.\allowbreak lean} has 5,611 lines, sha \texttt{b4d07119…}, its first 5,344
lines are byte-identical to \texttt{scratch/\allowbreak ChainEvenCorner.\allowbreak lean}, and its last 267 lines are byte-identical
to the promoted suffix; it compiles in \textbf{58 s on jack}, and its \texttt{\#print\ axioms} block names all four
new theorems. The promotion timings quoted in the round-26/26b notes (71 s jack, 74 s john) and the
john-side build, gate, tamper, axiom and scratch receipts (\textbf{67 s john}) are recorded in
\texttt{notes/\allowbreak math-papers.\allowbreak md} and were \textbf{not} re-verified for this draft: john was unreachable from
the evening of 2026-11-19 until 2026-11-25, his published battery report carries no round-26 timing,
and only the
statement layers were re-run on his machine since (\S{}5.5). They are labelled accordingly in \S{}5.5 and
Appendix A.

Two harness defects from the same overload event must be stated, because they explain two rows of the
table rather than a gap in the library. The report's \texttt{\#\#\ Layers\ 4,5,8} section is empty: the batch
driver died with a \texttt{UnicodeDecodeError} on a child's stdout cut mid-character, and the traceback sits
at lines 14--17 of the \texttt{.\allowbreak log}, so the \textbf{mutation/contradiction child's section is cut off there}; the
log itself continues to its last line and ends with \texttt{LEANCHECK\ DONE}, and the verdict table is correct
because it is computed from \texttt{results.\allowbreak jsonl}. And the scratch linter leg of layer 7 failed with
\texttt{FileNotFoundError:\ \textquotesingle{}lake\textquotesingle{}} because the subprocess did not inherit the elan PATH; it was re-run
correctly in layer 10 (\texttt{layer10/\allowbreak REPORT\_\allowbreak L10.\allowbreak md} step 4) and returned \textbf{0 warnings across 37 files}, of
which 31 compile --- the 6 failures are 3 stale library copies (\texttt{G.\allowbreak lean}, \texttt{General.\allowbreak lean}, \texttt{T.\allowbreak lean}) and
3 broken fragments (\texttt{E.\allowbreak lean}, a renamed Mathlib constant; \texttt{j28\_\allowbreak test.\allowbreak lean}, an unknown module prefix;
\texttt{j31\_\allowbreak tail.\allowbreak lean}, a fragment without imports).

One ordering fact a reader should have: the battery ran on 2026-11-19 from 02:07Z to 19:28Z, and
\texttt{Clifford/\allowbreak BridgeL10.\allowbreak lean} and the new root were written at 21:18Z the same day. Layers 2--9 therefore
never saw the bridge file; it is covered only by the build, gate, tamper and axiom outputs recorded in
\texttt{layer10/\allowbreak REPORT\_\allowbreak L10.\allowbreak md}, which post-date it.


\hypertarget{coverage-limits}{%
\section{Coverage limits}\label{coverage-limits}}

A formalisation paper is only as good as its account of what it has not done. These are the limits,
in descending order of consequence for a reader.

\begin{enumerate}
\def\labelenumi{\arabic{enumi}.}
\item
  \textbf{Mutation and contradiction probes now reach the late block, and neither covers every declaration.} The eight audited
  source files contain 577 \texttt{theorem}/\texttt{lemma} declarations. Mutation ran twice: on 2026-11-19 over
  259 declarations (198 outside \texttt{MorleyKernel.\allowbreak lean} and 61 of its 343, all in lines 33--806), and in
  a late block on 2026-11-24/25 over 257 further \texttt{MorleyKernel.\allowbreak lean} declarations, all in lines
  816--5586 and disjoint from the first set. \texttt{MorleyKernel.\allowbreak lean} therefore has \textbf{318 of its 343
  declarations (93\%)} under mutation, including the round 19--26 material: \texttt{heven\_\allowbreak kernel}
  (\href{https://github.com/jpbald93/paperl_clifford_chain/blob/d15464f0f4f4c8848e7d3e63dd68a03ce45433ae/Clifford/MorleyKernel.lean\#L4697}{\texttt{Clifford/\allowbreak MorleyKernel.lean}:4697}), \texttt{heven\_\allowbreak line\_\allowbreak of\_\allowbreak P0} (\href{https://github.com/jpbald93/paperl_clifford_chain/blob/d15464f0f4f4c8848e7d3e63dd68a03ce45433ae/Clifford/MorleyKernel.lean\#L5137}{\texttt{Clifford/\allowbreak MorleyKernel.lean}:5137}),
  \texttt{evenCorner\_\allowbreak of\_\allowbreak QL}, \texttt{heven\_\allowbreak inf\_\allowbreak all} and \texttt{polyF\_\allowbreak coeff\_\allowbreak top} (\href{https://github.com/jpbald93/paperl_clifford_chain/blob/d15464f0f4f4c8848e7d3e63dd68a03ce45433ae/Clifford/MorleyKernel.lean\#L5499}{\texttt{Clifford/\allowbreak MorleyKernel.lean}:5499})
  each carry mutant rows. The late block contributes 891 mutant rows --- 872 killed outright, 3 killed by heartbeat
  timeout, \textbf{16 survivors}, 0 unattributed --- against the first battery's 631 rows over all eight
  files. Of the sixteen late-block survivors, nine compile their own mutant proof and are caught
  downstream (at lines 1647, 2761, 3922 and 5574); three change \texttt{2*0+5} to \texttt{3*0+5} in a \texttt{k\ =\ 0}
  statement, where both are 5, so the mutant is the same proposition and the row is vacuous; and
  four alter the conjunct \texttt{u\ ≠\ 0} in the \emph{conclusion} of \texttt{child\_\allowbreak line\_\allowbreak of\_\allowbreak parent\_\allowbreak mid\_\allowbreak zero} --- a branch
  lemma whose hypotheses are inconsistent on purpose (\texttt{hmid} and \texttt{P₀\ ≠\ 0} force \texttt{V\ =\ 0}, contradicting
  \texttt{hv}), so no proof obligation reaches the mutant (\S{}5.8). \textbf{No survivor is a hole in a headline statement.} The 25 declarations with no mutant on
  jack split in two: six primed names that jack's collector skips
  (\texttt{d7\_children\_gcircleS\textquotesingle{}}/\texttt{\textquotesingle{}\textquotesingle{}}, \texttt{hodd\_inf\textquotesingle{}}, \texttt{hodd\_inf\_all\textquotesingle{}}, \texttt{hodd\_one\_P0\_line\textquotesingle{}},
  \texttt{p0\_levelpts\_collinear\textquotesingle{}}), which john's second-machine run mutates and kills every mutant of, and
  nineteen on which the generator produced no mutant --- \texttt{evenCornerQL\_\allowbreak proof}, \texttt{EvenCorner\_\allowbreak proof}
  and \texttt{cliffordChainS\_\allowbreak inf} are of that kind --- which are not covered by mutation: they rest on the
  gate, tamper and axiom checks and, where the declaration carries its hypotheses as explicit
  binders, on the satisfiability layer --- a hypothesis packed inside a \texttt{Prop} definition is invisible
  to both (\S{}5.2). Contradiction probes ran in two blocks as well: the battery's 294 declarations, all 234
  outside \texttt{MorleyKernel.\allowbreak lean} and 60 of its 343, every one before line 816, and on 2026-11-25 a late
  block over 257 further \texttt{MorleyKernel} declarations at lines 816--5592 --- 163 \texttt{ok}, 93 with no
  hypothesis and one \texttt{probe-invalid}, \textbf{0 CONTRADICTORY}. That is \textbf{551 of the 577 declarations
  (95.4\%)} probed. The 26 without a probe row are the six primed names jack's collector skips
  (\texttt{d7\_children\_gcircleS\textquotesingle{}}/\texttt{\textquotesingle{}\textquotesingle{}}, \texttt{hodd\_inf\textquotesingle{}}, \texttt{hodd\_inf\_all\textquotesingle{}}, \texttt{hodd\_one\_P0\_line\textquotesingle{}},
  \texttt{p0\_levelpts\_collinear\textquotesingle{}}) and twenty further \texttt{GenPosN}/determinant lemmas and definitional
  restatements (\texttt{cliffordChainS\_\allowbreak inf} --- the headline, not a corollary --- \texttt{cliffordChainS\_\allowbreak of\_\allowbreak corner}, \texttt{cliffordChainS\_\allowbreak of\_\allowbreak inf},
  \texttt{cliffordChainS\_of\_inf\textquotesingle{}}, \texttt{evenCornerQL\_\allowbreak proof}, \texttt{EvenCorner\_\allowbreak proof}, \texttt{evenCorner\_\allowbreak of\_\allowbreak cornerQL},
  \texttt{parent\_\allowbreak on\_\allowbreak child\_\allowbreak line}, \texttt{child\_\allowbreak det\_\allowbreak M7}, \texttt{child\_\allowbreak det\_\allowbreak MN}, \texttt{parent\_\allowbreak det\_\allowbreak MN}, \texttt{c5\_\allowbreak genPosN},
  \texttt{c6\_\allowbreak genPosN}, \texttt{d7\_\allowbreak genPosN}, \texttt{stepEven\_\allowbreak gen}, \texttt{stepOdd\_\allowbreak core}, \texttt{blockRev\_\allowbreak val}, \texttt{child\_\allowbreak line\_\allowbreak P0},
  \texttt{child\_\allowbreak line\_\allowbreak P0N}, \texttt{p0\_\allowbreak levelpts\_\allowbreak gcircleS}). One late-block row is invalid rather than passing: for
  \texttt{subfamily\_\allowbreak P0\_\allowbreak zero\_\allowbreak gen} the \texttt{True} control itself failed to compile (\texttt{conj} is bound as an
  implicit variable under \texttt{autoImplicit}), so the probe never ran and no contradiction claim is
  made for it. The 26 are a limit of jack's collector only: john's like-for-like receipt probes all
  343 \texttt{MorleyKernel} declarations, these 26 included, and reports 0 CONTRADICTORY
  (\texttt{layer0/\allowbreak john\_\allowbreak receipt\_\allowbreak 20261119/\allowbreak RECEIPT\_\allowbreak VERDICT.\allowbreak md}). \textbf{OPEN}: an expected-job manifest for the
  late block.
\item
  \textbf{Satisfiability was run on two machines; about a third of the declarations remain unresolved.} The
  2026-11-19 battery ran with \texttt{-\/\allowbreak -fast}, which replaces the job list with \texttt{sat\ =\ {[}{]}}: that row was a
  disabled layer, not an empty candidate class. The layer was then run properly on 2026-11-25/26 ---
  jack over 596 declarations (13 h 04 min) certifies \textbf{400 (67.1\%)}, john over 584 (10 h 17 min)
  certifies \textbf{394 (67.5\%)} --- and nothing was refuted: \texttt{stmt-invalid} is 0 on both machines, so no
  declaration's existential statement failed to elaborate. Vacuity is not inferred from that count.
  The 397 split by statement class: \textbf{210} have no hypothesis to inhabit, \textbf{119} have no binders
  at all (the statement reduced to \texttt{True}), and \textbf{68} inhabit a hypothesis of the declaration. So
  layer 8 witnesses the hypotheses of \textbf{68} of the 577 shared declarations, not 397 --- and a
  hypothesis packed inside a \texttt{Prop} definition (\texttt{EvenCornerQL}, \texttt{EvenCorner}, the payloads) is
  invisible to it, exactly as it is to layer 6. The one declaration whose hypotheses are the corner
  itself, \texttt{parent\_\allowbreak crossVec\_\allowbreak zero\_\allowbreak of\_\allowbreak mid}, is unresolved on both machines, so whether the corner is
  inhabited rests on the numerical witness of \S{}1.1, not on layer 8; for the remaining 180 nothing is
  claimed either way. Across
  the \textbf{577 declarations both machines ran, 560 agree} on the verdict class, 17 differ only in
  which machine found a witness, and \textbf{397 (68.8\%) are certified by at least one machine}; \textbf{180}
  are unresolved on both. A certificate is an accepted instance: Lean checks a proof of
  \texttt{∃\ binders,\ True} whose binders include each of the declaration's hypotheses as a \texttt{Prop}-valued
  binder, so the certificate must \emph{exhibit an inhabitant of every hypothesis}; that is what makes it
  a satisfiability witness rather than an assumption of what it is meant to certify. An \texttt{unresolved} row is a failed
  search --- deterministic patterns plus at most three model proposals --- and is \textbf{not} a refutation;
  it says nothing about the declaration one way or the other. A curated witness table (concrete data
  values supplied per declaration, 107 of the more tractable unresolved names, jack, 2026-11-26)
  closed \textbf{none}: 88 of the 107 fail because the instrument cannot discharge the propositional side
  conditions of the hypotheses --- its simp set leaves \texttt{(starRingEnd\ ℂ)\ 2} unreduced and \texttt{decide}
  cannot discharge equalities over \texttt{ℂ} --- so a table helps only together with a per-family proof
  library for \texttt{Transversal}, \texttt{GenTriple}, \texttt{GenPos4/\allowbreak 5}, \texttt{Chain\ 4\ !{[}\ldots{}{]}} and \texttt{mD\ ≠\ 0}. Layer 5's \texttt{ok}
  remains a failed search for \texttt{False}, not a satisfiability certificate, and layer 8b reports 0 links
  because neither tree contains a hand-written \texttt{*\_\allowbreak sat} witness. \textbf{OPEN}: the 180 unresolved
  declarations, with the per-family proof library the table would need.
\item
  \textbf{The layer-13 comparison covers all seven statements, and the bridge layer now closes two rows in
  both directions.} \texttt{L13\_\allowbreak FULL\_\allowbreak COMPARE.\allowbreak md} and \texttt{layer13/\allowbreak BlindCompareFull.\allowbreak lean} (compile exit 0,
  2.8 s, no \texttt{sorry}) compare the blind writer's real types --- \texttt{BlindGenPos}, \texttt{BlindSecond},
  \texttt{BlindIsGCircle} with \texttt{0\ \textless{}\ r}, \texttt{BlindCliffordP}, \texttt{BlindCliffordFin} --- against the paired library
  theorems. Row S6 is a \textbf{definitional} difference: the two sides assert the same geometric content over
  independently invented predicates, and the bridge proves it in both directions. S7 is definitional
  in its predicates but stronger in its hypothesis --- its writer again quantifies over \texttt{BlindGenPos},
  as S1 does --- and it is closed only because \texttt{N7\_\allowbreak RB} proves the writer's statement outright, without
  passing through the library theorem. Rows S4
  and S5 are not: the writer's finite-plane variant (\texttt{BlindCliffordFin}) requires every level object
  to lie in \texttt{BlindPlane}, so a level object that is a line --- which contains $\infty$ --- is unrepresentable,
  while the library's \texttt{CliffordChain} admits lines. That is a difference in the predicate, not in
  presentation: the \texttt{BlindGenPos} triple of \texttt{layer13/\allowbreak BRIDGE\_\allowbreak NOTES.\allowbreak md} (three centres on a circle
  through \texttt{M}, whose second intersections are distinct and collinear) satisfies the writer's
  hypotheses, and a blind finite-plane level-3 object would have to be a positive-radius circle
  through three distinct collinear points (impossible) or a line (excluded by the plane conjunct),
  so the writer's S4 statement is false as written at \texttt{n\ =\ 3}; the library's cutoff is unaffected
  and includes \texttt{n\ =\ 3}. This is a hand argument on exact coordinates --- the triple's second points are
  \texttt{2+2i}, \texttt{(1−√3/\allowbreak 2)+i(1+√3/\allowbreak 2)} and \texttt{3+i√3}, distinct and collinear --- not a Lean-checked negative
  statement. S5 is stated over the same encoding. Rows S2 and S3 add a
  quantifier difference: the writer asks for existence plus incidence of \textbf{every} certified point
  on one generalized circle, while the real payloads place the canonical \texttt{EptInf} only; the
  uniqueness step that would identify the two is false (\S{}5.4). S1 is definitional \emph{as well as}
  carrying a hypothesis-strength gap --- the writer quantifies over \texttt{BlindGenPos}, the library over the
  strictly stronger \texttt{GenPosN} --- so its two sides do \textbf{not} assert the same content, which is why it
  is excluded from the two definitional rows above and stays open. Because the real propositions are
  theorems, the blind statement implies each of them trivially --- that direction is closed for all
  seven rows and carries no information. The informative direction, real theorem $\Rightarrow$ writer's
  statement, is now \textbf{proved in Lean for S7 and S6}. \texttt{layer13/\allowbreak BridgeBlind.\allowbreak lean} (sha
  \texttt{f3819d69…}, 829 lines, compiles with exactly two \texttt{sorry}s, both at S1 and the general
  transport) supplies the predicate bridges --- the circle (\texttt{mem\_\allowbreak sphere\_\allowbreak iff\_\allowbreak onCircleM}), the
  transversal/second-point pair (\texttt{blindTransversal\_\allowbreak iff}, \texttt{BlindSecond\_\allowbreak eq}), general position
  (\texttt{genPosN\_\allowbreak blindGenPos}), the generalized-circle transport (\texttt{isGCircleS\_\allowbreak of\_\allowbreak blindIsGCircle},
  \texttt{chainS\_\allowbreak of\_\allowbreak blind}), and \texttt{BlindSix\_\allowbreak eq} --- and on them \texttt{N7\_\allowbreak RB} (proved unconditionally: the writer's
  level-3 statement holds outright) and \texttt{N6\_\allowbreak RB} (existence by an explicit per-level construction
  over the six \texttt{Counter6} centres, uniqueness by \texttt{chainS\_\allowbreak six\_\allowbreak unique}). Both depend only on
  \textbf{std}, with no \texttt{sorry} in their cone. \textbf{S1 stays open and is not a bridging problem}: \texttt{RealN1}
  supplies chains only for \texttt{GenPosN}, the writer quantifies over the strictly weaker \texttt{BlindGenPos},
  and a concrete three-centre configuration separates the two (\texttt{layer13/\allowbreak BRIDGE\_\allowbreak NOTES.\allowbreak md}); whether
  S1 holds at the weaker hypothesis is unsettled, and settling it needs a direct writer-side
  construction, not a rescoping of the statement: \texttt{layer13/\allowbreak BRIDGE\_\allowbreak NOTES.\allowbreak md} \S{}4 records that the
  blind statement appears to be true at the weaker hypothesis (the counterexample search above found
  none), so weakening it to \texttt{GenPosN} would replace a true statement by a weaker one. That
  construction is partly recorded: \texttt{layer13/\allowbreak BlindChain.\allowbreak lean} (1,296 lines, sha \texttt{6d5bfe5d…}) proves
  the writer's level-3 and level-4 objects for every \texttt{BlindGenPos} tuple, including the collinear and
  mixed-degenerate cases, proves level 5 under \texttt{GenPos5} or finite Miquel points, and reduces the
  all-\texttt{n} statement to one named lemma, \texttt{blindClifford\_\allowbreak second} --- the all-degenerate deletion --- which
  is the file's single remaining \texttt{sorry}.
  A search for a counterexample to the blind statement found none
  (\texttt{layer13/\allowbreak s1\_\allowbreak counterexample/\allowbreak SUMMARY.\allowbreak md}): all 3,732 and 25,820 \texttt{BlindGenPos} configurations on a
  small integer grid at \texttt{n\ =\ 4,\ 5}, 1,197 further random configurations checked at sixty digits, and
  a sample at \texttt{n\ =\ 6}, all admit the blind object --- the collinear case is absorbed because a blind
  line carries \texttt{∞}. So S1 is expected to be true as the writer states it, but it is not proved, and
  this paper does not claim it.
  The general transport \texttt{blindOfChainS} also stays open on one branch: \texttt{IsGCircleS} admits the
  radius-zero circle (a finite singleton) where \texttt{BlindIsGCircle} requires \texttt{0\ \textless{}\ r}. \textbf{OPEN}: S2/S3
  need a universal-incidence theorem, not uniqueness; S4/S5 need a corrected writer-side encoding of
  the finite plane (S4 is false as written, because the encoding forbids the line case that the
  library admits); S1 needs the statement decision above. The comparison file itself does run on a
  second machine: \texttt{BlindCompareFull.\allowbreak lean} (sha \texttt{32dfa1b2…}) is byte-identical to
  \texttt{john\allowbreak :\textasciitilde/\allowbreak lean/\allowbreak paperL\_\allowbreak john/\allowbreak layer13\_\allowbreak receipt.\allowbreak lean} and compiles there (rc 0, 5.4 s against jack's
  2.8 s); the blind writer was not re-run on john.
\item
  \textbf{The even-level Clifford point is not unique, and \texttt{UniqEven} is false.} \texttt{EptInf} is a canonical
  selection, not a theorem of uniqueness: for the centres \texttt{2t\ −\ 2it²} the level-6 objects \texttt{\{4i\}}
  and \texttt{\{$\infty$\}} both satisfy \texttt{ChainS\ 6}, as checked in \texttt{check/\allowbreak not\_\allowbreak UniqEven.\allowbreak lean} (\textbf{std} axioms). Only
  the \texttt{Counter6} instance has a uniqueness theorem (\texttt{chainS\_\allowbreak six\_\allowbreak unique}). Two consequences: the
  definite article of the former \S{}2.5 wording (now replaced by ``some level-(n$-$1) object'') cannot be
  justified in general, and any future statement
  identifying existential witnesses with \texttt{EptInf} needs extra hypotheses or a different selection
  predicate. Nothing here weakens the existence theorem.
\item
  \textbf{The two-machine receipt now holds for the statement layers.} The series' evidence standard is
  a second machine. On 2026-11-24/25 john audited the same file at the same pin
  (\texttt{Clifford/\allowbreak MorleyKernel.\allowbreak lean}, sha \texttt{dce7b011…}) with the same generator under a patched \texttt{REPOS}
  and \texttt{OUT}: \textbf{1,083 mutants over 324 declarations --- 1,058 killed, 9 killed by heartbeat timeout,
  16 SURVIVED --- and 343 contradiction probes (228 \texttt{ok}, 114 no-hypothesis, 1 probe-invalid, 0
  timeout), 0 CONTRADICTORY.} That reproduces jack exactly where reproduction is possible: the
  survivor set is the \textbf{same sixteen \texttt{(declaration,\ mutation)} pairs} on both machines, no
  survivor is machine-specific, and both report zero contradictions (\S{}5.1). john also covers
  the six primed declarations jack's collector skips, killing all of their mutants. The v0.4 numbers
  quoted from \texttt{notes/\allowbreak math-papers.\allowbreak md} --- 1,536 mutant rows over 520 declarations with 17 survivors,
  and 573 contradiction declarations --- came from john's earlier wide run over a different tree and
  are \textbf{superseded} by this like-for-like rerun. What remains external: the round-26b timings (71 s on jack and 74 s on
  john for the promotion build, 67 s for the scratch receipt on john) and the build/gate/tamper receipts, quoted from the notes and not re-run
  here; and layers 6 and 9 were never run on john. Full comparison:
  \texttt{layer0/\allowbreak john\_\allowbreak receipt\_\allowbreak 20261119/\allowbreak RECEIPT\_\allowbreak VERDICT.\allowbreak md}.
\item
  \textbf{The ``75 unbridged notions'' is inventory arithmetic, not a semantic count.} Layer 9's 137 / 24 /
  113 comes from a heuristic \texttt{has-bridge} flag that points at internal lemmas as well as at
  Mathlib: its \texttt{has-bridge} list contains \texttt{GenPosN\ →\ hodd\_\allowbreak inf\_\allowbreak P0\_\allowbreak eq}, \texttt{KP\ →\ A7\_\allowbreak 0\_\allowbreak eq},
  \texttt{IsGCircleS\ →\ IsGCircleS.\allowbreak infty\_\allowbreak mem} and \texttt{ChainS\ →\ chainS\_\allowbreak six\_\allowbreak iff}, none of which identifies the
  notion with a Mathlib counterpart, while \texttt{GenTriple} stays among the 75 even though
  \texttt{genTriple\_\allowbreak iff} (\href{https://github.com/jpbald93/paperl_clifford_chain/blob/d15464f0f4f4c8848e7d3e63dd68a03ce45433ae/Clifford/Bridge.lean\#L110}{\texttt{Clifford/\allowbreak Bridge.lean}:110}) already bridges it. The arithmetic 113 $-$ 13 $-$ 25 = 75
  is right; the reading ``75 notions have no bridge to Mathlib'' is not. The reason the sphere
  machinery is unlikely to bridge is real and separate: Mathlib has no generalized-circle or
  M\"obius-sphere notion to target, so \texttt{IsGCircleS}, \texttt{ChainS}, \texttt{CliffordChainS} and \texttt{liftS} have
  nothing to be identified with. The 19 concrete bridge theorems and their 13 targets stand on their
  own. \textbf{OPEN}: a manual reconciliation of false positives and negatives.
\item
  \textbf{The 16 unattributed rows are now decisive.} Nine mutation rows and seven contradiction rows had
  carried \texttt{rc=137} (SIGKILL) or \texttt{rc=None} with an empty error string: the child was killed while the
  host sat at load $\approx$ 50. They were re-run on 2026-11-25, 19:32Z--19:42Z, under a lighter load and
  restricted to the 14 affected declarations (\texttt{lc\_\allowbreak mod\_\allowbreak requeue.\allowbreak py}, 46 jobs), and all sixteen
  resolved --- the nine mutation rows as kills, the seven probe rows as \texttt{ok}. The battery file now
  contains \textbf{935 rows and no \texttt{error-unattributed} row}: 631 mutation rows (535 killed, 95 killed by
  heartbeat or wall-clock timeout, 1 survivor) and 294 probe rows (147 \texttt{ok}, 125 no-hypothesis,
  22 timeout). The only rows still without a verdict are those 22 probe timeouts, which are limits
  of the tactic budget, not load artifacts, and the one \texttt{probe-invalid} row of item 1.
\item
  \textbf{Eight unused hypotheses on the layer-6 list, plus one the ablation cannot see.} \texttt{cramer2\_\allowbreak P0}
  and \texttt{cramer2\_\allowbreak Qm} both still take \texttt{(hm\ \allowbreak :\ 1\ ≤\ m)}; \texttt{child\_\allowbreak levelPt\_\allowbreak kernel\_\allowbreak line} takes
  \texttt{(hd\ \allowbreak :\ ∀\ j,\ d\ j\ ≠\ 0)}; \texttt{p0\_\allowbreak children\_\allowbreak gcircleS\_\allowbreak genPos} takes a \texttt{\_\allowbreak hP0} it does not use;
  \texttt{kpe\_\allowbreak odd\_\allowbreak succ\_\allowbreak eq\_\allowbreak mul\_\allowbreak x}, \texttt{kpe5\_\allowbreak hi}, \texttt{kpe6\_\allowbreak hi} take bound hypotheses \texttt{hj}, \texttt{hj\textquotesingle{}} that are unused;
  and \texttt{parent\_\allowbreak on\_\allowbreak child\_\allowbreak line} takes \texttt{(hQ0\ \allowbreak :\ …\ ≠\ 0)}. The late-block mutation run adds a different case, not an unused input:
  \texttt{child\_\allowbreak line\_\allowbreak of\_\allowbreak parent\_\allowbreak mid\_\allowbreak zero} concludes \texttt{∃\ p\ u,\ u\ ≠\ 0\ ∧\ …} from \texttt{hv\ \allowbreak :\ V\ ≠\ 0}, \texttt{hmid} and
  \texttt{P₀\ ≠\ 0}, and its proof is \texttt{absurd\ (parent\_\allowbreak crossVec\_\allowbreak zero\_\allowbreak of\_\allowbreak mid\ …)\ hv} --- the branch is empty, so
  the four mutants of the conjunct \texttt{u\ ≠\ 0} in its \emph{conclusion} have no proof obligation to break and
  survive. The lemma is an empty-branch closer that mutation cannot test; that is a limit of the
  audit, not an unused hypothesis, and it is why \texttt{0\ CONTRADICTORY} in layer 5 does not by itself
  exclude this source-visible vacuity. An unused hypothesis is not an error, but a
  reader should know that these eight \emph{auxiliary} lemmas carry assumptions their proofs do not use;
  removing them would strengthen those lemmas, and the headline statements of \S{}2 are unaffected
  either way. To them add \texttt{evenCornerQL\_\allowbreak proof}'s unused \texttt{\_\allowbreak hall} (\S{}2.4), which layer 6 could not
  reach because \texttt{EvenCornerQL} is a \texttt{Prop} definition. \textbf{OPEN}: tighten or document each.
\item
  \textbf{Carver's classification is cited, not formalised.} \S{}1.2 says this; it belongs in the limits
  list too. The development proves the chain exists on the sphere, including the line and $\infty$ cases
  as they arise; it does not state or prove the key-matrix criterion that predicts \emph{when} they
  arise, and its carrier is not Carver's (which admits a point-circle at infinity that
  \texttt{IsGCircleS} excludes). \textbf{OPEN}.
\item
  \textbf{The prior-art search has recorded gaps}, and a novelty claim is only as strong as its weakest
  search: MathSciNet and zbMATH were not searched (no subscription or API key available); Morley
  1900, F. Morley 1929, Wagner 1929 and Kobayashi 1958 were seen only as metadata behind the JSTOR
  paywall (Carver 1920 is the only one of them read in full, from the archive.org OCR); HAL timed
  out twice; Google Scholar is not reachable from this runtime; the non-English web query returned
  nothing; the Mizar MML body was not full-text searched; and LeanGeo, LeanPool, Principia, TLMC,
  ACL2, PVS, Nuprl and Naproche were not searched individually. \texttt{SCOUT.\allowbreak md} \S{}4 records these as
  \emph{not searched}, not as \emph{empty}. The novelty claim in \S{}1 is therefore: no formal statement of the
  chain for arbitrary n was found in any of the sources logged in \texttt{SCOUT.\allowbreak md} \S{}0 (26 search rows), not: none exists.
\item
  \textbf{What has and has not run, per input version, and on how many machines.} Layers 1--10 ran
  against the code (the battery at \texttt{dce7b011…}, before the bridge file; layer 10 after it). Layer 11
  ran on v0.1 and was \textbf{re-run} on every draft from v0.2 through v0.11; each re-run is a corrected
  re-run on a new input, not a reproduction, and v0.5--v0.11 changed only evidence and provenance, so the
  citation set it certifies is the one introduced between v0.2 and v0.4. Layer 12 \textbf{reviewed v0.1};
  the three-model round reviewed v0.2; an independent review of v0.3 is applied in v0.4; an
  independent referee pass on v0.7 produced the corrections of v0.8, a referee pass on v0.9 those
  of v0.10, and an Opus 5.5 pass on v0.10 those of v0.11 (Appendix C). Layer 13's blind
  writer was \textbf{re-run} after the parity repair of the blind text, three rows were compared on v0.2,
  and all seven rows were compared against the writer's real types on v0.3, with the bridge layer
  closing S6 and S7 in both directions (\S{}5.3). Layer 14 was cut from v0.2 and \textbf{re-cut and re-run on
  v0.7}; layer 15 was built on v0.2 and \textbf{re-built on v0.7} (\texttt{L14\_\allowbreak RESTATEMENTS\_\allowbreak v07.\allowbreak md},
  \texttt{L15\_\allowbreak PRIOR\_\allowbreak ART\_\allowbreak v07.\allowbreak md}).
  On machine count: layers 4 and 5 were reproduced on \textbf{both} machines for \texttt{MorleyKernel.\allowbreak lean} (the
  like-for-like receipt, \S{}5.5); layer 8 ran on \textbf{both} over their full inventories (\S{}5.2); and the
  layer-13 comparison file \texttt{BlindCompareFull.\allowbreak lean} compiles on both, while the later bridge file
  \texttt{BridgeBlind.\allowbreak lean} was checked on jack only. Layers 1, 2 and 7 also have a john run in his 19 November battery at the
  same pin (wider file set: gate PASS 36, tamper 8/8, linter 0; published in
  \texttt{layer0/\allowbreak john\_\allowbreak receipt\_\allowbreak 20261119/\allowbreak BATTERY.\allowbreak md}); layers 6, 9, 10, 11, 12, 14 and 15 are jack-only. The upload is no
  longer open: the pin is commit \texttt{d15464f}, the repository is public at that commit and the
  supporting material is published (\S{}3); the Zenodo deposit is the authors' remaining step and is
  not described here as done.
\end{enumerate}


\hypertarget{appendix-a-provenance-of-this-draft}{%
\section*{Appendix A --- provenance of this draft}\label{appendix-a-provenance-of-this-draft}}

All paths in this appendix are relative to the paper's artifact directory \texttt{AFM/\allowbreak PaperL/\allowbreak } unless
prefixed \texttt{jack\allowbreak :} or \texttt{john\allowbreak :} (host paths). A flat review packet that copies these files to one
directory must be read through this mapping; the files themselves are the ones named here.

\begin{itemize}
\tightlist
\item
  Claim set, route and statement form: \texttt{AFM/\allowbreak PaperL/\allowbreak PLAN.\allowbreak md} \S{}Claim set, \S{}Route.
\item
  Novelty verdicts quoted in \S{}1 and \S{}5.10: \texttt{AFM/\allowbreak PaperL/\allowbreak SCOUT.\allowbreak md} \S{}0 (search log, 26 rows), \S{}1
  (verdict table, L1--L9), \S{}4 (gaps), \S{}5 (what may and may not be claimed).
\item
  Rounds 19--26b (odd side, even side, the corner, \texttt{EvenCornerQL}, promotion, gate 32 $\to$ 36):
  \texttt{notes/\allowbreak math-papers.\allowbreak md}, sections ``Paper L --- round 19'' through ``Paper L --- round 26b''. Bundled
  excerpt: \texttt{evidence/\allowbreak math-papers-PaperL-excerpt.\allowbreak md} (also at \texttt{jack\allowbreak :\textasciitilde/\allowbreak paperL\_\allowbreak jack/\allowbreak layer0/\allowbreak }); it is the
  source of the john timings quoted in \S{}4 and \S{}5.5, which remain \emph{recorded, not verified}. Its john
  battery snapshot is superseded by the like-for-like rerun of 2026-11-25 (\S{}5.5).
\item
  Layer 10 content and the 19 bridges: \texttt{jack\allowbreak :\textasciitilde/\allowbreak paperL\_\allowbreak jack/\allowbreak layer10/\allowbreak REPORT\_\allowbreak L10.\allowbreak md}.
\item
  Layer 11 anchor table (36 declarations, 19 bridges, axioms): \texttt{jack\allowbreak :\textasciitilde/\allowbreak paperL\_\allowbreak jack/\allowbreak layer11/\allowbreak ANCHORS.\allowbreak md};
  the citation gate: \texttt{namegate.\allowbreak py} and \texttt{L11\_\allowbreak NAMEGATE.\allowbreak md} (regenerated for every draft from v0.4
  through v0.11; the current report is v0.11).
\item
  Layers 1--9 numbers, harness defects, the 16 unattributed rows, the layer-6 unused list, john's
  snapshot: \texttt{jack\allowbreak :\textasciitilde/\allowbreak leancheck/\allowbreak reports/\allowbreak paperL\_\allowbreak jack\_\allowbreak 20261119\_\allowbreak 020713.\allowbreak md} and
  \texttt{notes/\allowbreak math-papers.\allowbreak md} ``Paper L --- round 26b leancheck battery''.
\item
  Corner witness (numerical, not a certificate) and its margins:
  \texttt{layer0/\allowbreak r25\_\allowbreak jack/\allowbreak RANK\_\allowbreak ROUTE\_\allowbreak COUNTEREXAMPLE.\allowbreak json}, bundled with the paper.
\item
  The two check files added for this draft: \texttt{check/\allowbreak allN.\allowbreak lean} (the all-n composition, rc 0, \textbf{std})
  and \texttt{check/\allowbreak not\_\allowbreak UniqEven.\allowbreak lean} (the counterexample to uniqueness, rc 0, \textbf{std}); copies at
  \texttt{jack\allowbreak :\textasciitilde/\allowbreak paperL\_\allowbreak jack/\allowbreak scratch/\allowbreak }.
\item
  Paper D (base cases \texttt{n\ ≤\ 5} and the six-centre failure): Zenodo record 23653309.
\item
  John's 19 November battery report is published in \texttt{layer0/\allowbreak john\_\allowbreak receipt\_\allowbreak 20261119/\allowbreak } (\texttt{BATTERY.\allowbreak md},
  \texttt{.\allowbreak log}, \texttt{RESULTS\_\allowbreak 20261119.\allowbreak jsonl}); it was hosted only at
  \texttt{john\allowbreak :\textasciitilde/\allowbreak leancheck/\allowbreak reports/\allowbreak lean/\allowbreak paperL\_\allowbreak john\_\allowbreak 20261119\_\allowbreak 143939.\allowbreak md} before this release. Still host-only:
  jack's \texttt{leancheck} battery report and \texttt{.\allowbreak log}, the layer-6/7/9 raw outputs behind 501/8 and
  137/24/113 (\texttt{\textasciitilde/\allowbreak leancheck/\allowbreak reports/\allowbreak paperL\_\allowbreak jack\_\allowbreak 20261119\_\allowbreak 020713.\allowbreak md}, \texttt{\textasciitilde/\allowbreak lc\_\allowbreak gap\_\allowbreak paperL\_\allowbreak jack/\allowbreak }), and the
  round-26/26b timings for john. Those figures are quoted from the notes excerpt.
\end{itemize}

\hypertarget{where-each-cited-artifact-is-published}{%
\subsection{Where each cited artifact is published}\label{where-each-cited-artifact-is-published}}

{\scriptsize\begin{longtable}[]{@{}
  >{\raggedright\arraybackslash}p{(\columnwidth - 2\tabcolsep) * \real{0.5000}}
  >{\raggedright\arraybackslash}p{(\columnwidth - 2\tabcolsep) * \real{0.5000}}@{}}
\toprule\noalign{}
\begin{minipage}[b]{\linewidth}\raggedright
cited here
\end{minipage} & \begin{minipage}[b]{\linewidth}\raggedright
published in \texttt{LEAN-verification/\allowbreak PaperL}
\end{minipage} \\
\midrule\noalign{}
\endhead
\bottomrule\noalign{}
\endlastfoot
\texttt{AFM/\allowbreak PaperL/\allowbreak DRAFT.\allowbreak md} & \texttt{preprint/\allowbreak DRAFT\_\allowbreak v11.\allowbreak md} and the PDF beside it \\
\texttt{AFM/\allowbreak PaperL/\allowbreak PLAN.\allowbreak md}, \texttt{SCOUT.\allowbreak md}, \texttt{namegate.\allowbreak py}, \texttt{L11\_\allowbreak NAMEGATE.\allowbreak md}, \texttt{L12\_\allowbreak REVIEW.\allowbreak md}, \texttt{L13\_\allowbreak REPORT.\allowbreak md}, \texttt{L14\_\allowbreak *}, \texttt{L15\_\allowbreak *}, \texttt{JOHN\_\allowbreak LEG.\allowbreak md} & \texttt{reports/\allowbreak } \\
\texttt{AFM/\allowbreak PaperL/\allowbreak REVIEW\_\allowbreak *}, \texttt{REVIEW\_\allowbreak TRIAGE\_\allowbreak *} & \texttt{reviews/\allowbreak } \\
\texttt{AFM/\allowbreak PaperL/\allowbreak layer13/\allowbreak *} & \texttt{layer13/\allowbreak } (the same files, minus \texttt{run4/\allowbreak }--\texttt{run6/\allowbreak }) \\
\texttt{AFM/\allowbreak PaperL/\allowbreak layer0/\allowbreak *} (the eight run directories cited) & \texttt{layer0/\allowbreak } \\
\texttt{AFM/\allowbreak PaperL/\allowbreak check/\allowbreak *}, \texttt{evidence/\allowbreak *} & \texttt{check/\allowbreak }, \texttt{evidence/\allowbreak } \\
\texttt{layer13/\allowbreak BlindCompareFull.\allowbreak lean}, \texttt{layer13/\allowbreak BridgeBlind.\allowbreak lean} & \texttt{layer13/\allowbreak } in this folder; in the code repository they are \texttt{scratch/\allowbreak …} \\
\end{longtable}}

\hypertarget{appendix-b-which-layers-ran-and-where-their-output-lives}{%
\section*{Appendix B --- which layers ran, and where their output lives}\label{appendix-b-which-layers-ran-and-where-their-output-lives}}

{\scriptsize\begin{longtable}[]{@{}
  >{\raggedright\arraybackslash}p{(\columnwidth - 6\tabcolsep) * \real{0.2500}}
  >{\raggedright\arraybackslash}p{(\columnwidth - 6\tabcolsep) * \real{0.2500}}
  >{\raggedright\arraybackslash}p{(\columnwidth - 6\tabcolsep) * \real{0.2500}}
  >{\raggedright\arraybackslash}p{(\columnwidth - 6\tabcolsep) * \real{0.2500}}@{}}
\toprule\noalign{}
\begin{minipage}[b]{\linewidth}\raggedright
layer
\end{minipage} & \begin{minipage}[b]{\linewidth}\raggedright
input
\end{minipage} & \begin{minipage}[b]{\linewidth}\raggedright
what it establishes
\end{minipage} & \begin{minipage}[b]{\linewidth}\raggedright
artifact
\end{minipage} \\
\midrule\noalign{}
\endhead
\bottomrule\noalign{}
\endlastfoot
1 & code \texttt{dce7b011…} & gate: 36 REQUIRED declarations present, axioms $\subseteq$ std & \texttt{gate.\allowbreak sh}, \S{}4 \\
2 & code \texttt{dce7b011…} & tampering: 8 cheat attempts rejected & \texttt{tests/\allowbreak tamper.\allowbreak sh} \\
baseline & code \texttt{dce7b011…} & the 8 audited source files compile & battery report \\
4 & code \texttt{dce7b011…} & mutation: 535 decisive kills, 1 survivor, 0 unattributed over 259 declarations --- plus a late block of 891 rows (872 kills, 3 timeout, 16 survivors) over 257 further \texttt{MorleyKernel} declarations; 318/343 of that file in total (\S{}5.1) & battery report; \texttt{layer0/\allowbreak lateblock\_\allowbreak jack/\allowbreak } \\
5 & code \texttt{dce7b011…} & contradiction probes: 0 CONTRADICTORY over 551 declarations --- 294 from the battery and 257 from the late block at lines 816--5592 (failed proof search, \S{}5.2) & battery report; \texttt{layer0/\allowbreak contra\_\allowbreak late\_\allowbreak jack/\allowbreak } \\
6 & code \texttt{dce7b011…} & ablation: 501 needed / 8 unused over 262 declarations & battery report, \S{}5.8 \\
7 & code \texttt{dce7b011…} & linter: 0 warnings in 8 library files; scratch leg from layer 10: 0 in 37 files, 31 compiling & battery report; \texttt{layer10/\allowbreak REPORT\_\allowbreak L10.\allowbreak md} \\
8 & code \texttt{dce7b011…} & satisfiability certification: \textbf{run on jack and john}, 400/596 and 394/584 certified; 397 of 577 shared declarations certified by at least one machine, 180 by neither & \texttt{layer0/\allowbreak sat\_\allowbreak jack/\allowbreak }, \texttt{layer0/\allowbreak sat\_\allowbreak john/\allowbreak }, \S{}5.2 \\
8b & code \texttt{dce7b011…} & witness linkage: \textbf{0 links} (no \texttt{*\_\allowbreak sat} witnesses in either tree) & \S{}5.2 \\
9 & code \texttt{dce7b011…} & notion inventory: 137 / 24 flagged / 113, heuristic & \texttt{\textasciitilde/\allowbreak lc\_\allowbreak gap\_\allowbreak paperL\_\allowbreak jack/\allowbreak inventory.\allowbreak md} \\
10 & code + \texttt{BridgeL10.\allowbreak lean} & definition bridges: 19 theorems, 13 notions bridged, 25 INTERNAL, 75 left & \texttt{layer10/\allowbreak REPORT\_\allowbreak L10.\allowbreak md} \\
11 & v0.1, re-run on v0.2 through v0.11 (evidence-only edits in v0.5--v0.11) & citation gate: 143 occurrences / 98 distinct, all resolving; 85 table rows, 30 of them signature anchors & \texttt{L11\_\allowbreak NAMEGATE.\allowbreak md} \\
12 & v0.1, v0.2 (three models), v0.3, v0.7, v0.9, v0.10 & prose versus code (independent reviewers): v0.1 review 5 MAJOR / 6 MINOR $\to$ v0.2; three-model round on v0.2 $\to$ 19 rows $\to$ v0.3; v0.3 review 0 MAJOR / 6 MINOR / 3 NOTE $\to$ v0.4; v0.7 referee 7 MAJOR / 3 MINOR $\to$ v0.8; v0.9 referee 2 MAJOR / 8 MINOR $\to$ v0.10; v0.10 Opus 5.5 pass $\to$ v0.11 & \texttt{L12\_\allowbreak REVIEW.\allowbreak md}, \texttt{REVIEW\_\allowbreak *\_\allowbreak v02.\allowbreak md}, \texttt{REVIEW\_\allowbreak TRIAGE\_\allowbreak v02.\allowbreak md}, \texttt{REVIEW\_\allowbreak gpt-6-astra\_\allowbreak v03.\allowbreak md}, \texttt{REVIEW\_\allowbreak TRIAGE\_\allowbreak v03b.\allowbreak md} \\
13 & v0.3 prose, blind writer re-run after the parity repair & blind re-formalisation: seven rows compared against the writer's real types; the bridge layer now proves the real $\Rightarrow$ writer direction for S6 and S7 (\S{}5.3); S4/S5 are stated over the writer's finite-plane encoding, which excludes lines, and S4 is false as written there & \texttt{layer13/\allowbreak L13\_\allowbreak REPORT.\allowbreak md} and \texttt{layer13/\allowbreak L13\_\allowbreak FULL\_\allowbreak COMPARE.\allowbreak md}, with \texttt{jack\allowbreak :\textasciitilde/\allowbreak paperL\_\allowbreak jack/\allowbreak scratch/\allowbreak } for the comparison sources and \texttt{jack\allowbreak :\textasciitilde/\allowbreak reformal/\allowbreak } for the blind files \\
14 & v0.2, re-cut and re-run on v0.7 & independent restatement of the seven signatures: 7 of 7 faithful; two wording notes & \texttt{L14\_\allowbreak RESTATEMENTS\_\allowbreak v07.\allowbreak md} (v0.2 pass: \texttt{layer13/\allowbreak L14\_\allowbreak RESTATEMENTS.\allowbreak md}) \\
15 & v0.2 + \texttt{SCOUT.\allowbreak md}, rebuilt on v0.7 & prior-art verdict table: PASS with recorded search limits & \texttt{L15\_\allowbreak PRIOR\_\allowbreak ART\_\allowbreak v07.\allowbreak md} (v0.2 table: \texttt{L15\_\allowbreak PRIOR\_\allowbreak ART.\allowbreak md}) \\
\end{longtable}}

Second-machine reproduction: layers 4 and 5 are reproduced on john at the same pin for
\texttt{MorleyKernel.\allowbreak lean} --- 1,083 mutants over 324 declarations, the identical sixteen survivors, zero
contradictions (\texttt{layer0/\allowbreak john\_\allowbreak receipt\_\allowbreak 20261119/\allowbreak RECEIPT\_\allowbreak VERDICT.\allowbreak md}) --- and the layer-13 comparison file
\texttt{BlindCompareFull.\allowbreak lean} also compiles there byte-identically; the later bridge file was compiled on
jack only. Layers 1, 2 and 7 also ran on john in his 19 November battery, on a wider file set; layers 6 and 9 were not run on john; the round-26b build/gate/tamper timings
are still quoted from the notes; and the re-runs of layers 11 and 13 are re-runs on new inputs, not
reproductions.

\hypertarget{appendix-c-changes-in-v0.3-through-v0.12}{%
\section*{Appendix C --- changes in v0.3 through v0.12}\label{appendix-c-changes-in-v0.3-through-v0.12}}

\hypertarget{v0.3-all-nineteen-rows-of-review_triage_v02.md}{%
\subsection{\texorpdfstring{v0.3 --- all nineteen rows of \texttt{REVIEW\_\allowbreak TRIAGE\_\allowbreak v02.\allowbreak md}}{v0.3 --- all nineteen rows of REVIEW\_TRIAGE\_v02.md}}\label{v0.3-all-nineteen-rows-of-review_triage_v02.md}}

All nineteen rows are applied; none required a change to the Lean.

\begin{enumerate}
\def\labelenumi{\arabic{enumi}.}
\tightlist
\item
  \S{}2.4/\S{}5.4: \texttt{UniqEven} is \textbf{false}, with the checked counterexample; the named-point form is a
  selection, not a uniqueness theorem.
\item
  \S{}5.3: the layer-13 comparison is described as it was --- writer's universal reading, re-typed
  existential restatement, three rows, one direction; ``structure untouched'' and ``strictly stronger''
  withdrawn. \texttt{L13\_\allowbreak REPORT.\allowbreak md} corrected to match. (v0.3$\to$v0.4: rewritten again for the seven-row
  comparison of \S{}5.3 below.)
\item
  \S{}4/\S{}5.2: layer 8 is disabled by \texttt{-\/\allowbreak -fast}; layer 8b has no recorded result; layer 5's \texttt{ok} is
  failed proof search.
\item
  \S{}4/\S{}5.1: mutation/contradiction coverage disclosed --- 61/60 of 343, all before line 807, no headline
  theorem --- and separated from the 16 inconclusive rows.
\item
  \S{}4/\S{}5.5/App. A: john's receipts and timings labelled recorded-not-verified; the corner witness
  bundled and labelled numerical; grandchildren \textbf{coincide} in the exact corner.
\item
  \S{}3.3: signature anchors added for the statement declarations; the gate claim now matches what the
  script checks (occurrence count, comment and comma forms, no bare ``line N'').
\item
  \S{}2.1: \texttt{second}, \texttt{Transversal}, \texttt{GenTriple} (both conjuncts), \texttt{crossVec}, \texttt{KP}, \texttt{levelPointN} and
  \texttt{EptInf} (with its zero-vector totalization) defined; \S{}2.3 adds the role of \texttt{h4\_\allowbreak pt\_\allowbreak inf} and
  \texttt{cliffordChainS\_\allowbreak all\_\allowbreak of}.
\item
  \S{}1.2: Paper D's C$_6$ qualified; base cases limited to \texttt{n\ ≤\ 5}; \texttt{not\_\allowbreak cliffordChain\_\allowbreak of\_\allowbreak six\_\allowbreak le} named
  for the \texttt{n\ ≥\ 6} failure.
\item
  \S{}1.1/\S{}2.5: \texttt{cliffordChainS\_\allowbreak of\_\allowbreak le\_\allowbreak five} cited for the base cases; the aggregate all-n statement is
  the composition in \texttt{check/\allowbreak allN.\allowbreak lean}, since no single declaration states it.
\item
  \S{}1/\S{}1.2/\S{}2.2: Carver cited as \S{}5 pp.~160--161 (Theorem 13) and \S{}6 pp.~161--164 (limiting
  figures); the point-circle-at-infinity difference from \texttt{IsGCircleS} stated.
\item
  \S{}4/\S{}5.6: layer 9's 24/113/75 relabelled as inventory flags, with the false positive and false
  negative named.
\item
  \S{}4: the survivor is the alternative-level mutant \texttt{CliffordChain\ 1\ →\ 2}, caught downstream by a
  type mismatch; ``weaker''/``not counted'' removed.
\item
  \S{}4: 13 macro use sites (15 identifier occurrences), three reworded module comments.
\item
  \S{}5.11/App. B: per-layer input versions; the blanket ``nothing was run twice'' removed.
\item
  \S{}2.5: indefinite article throughout, with the reason.
\item
  App. A/App. B: artifact paths corrected --- paper-relative paths under the \texttt{AFM/\allowbreak PaperL/\allowbreak } tree,
  host paths for \texttt{jack\allowbreak :\textasciitilde/\allowbreak paperL\_\allowbreak jack/\allowbreak scratch/\allowbreak } and \texttt{jack\allowbreak :\textasciitilde/\allowbreak reformal/\allowbreak }; a flat review packet is
  read through the mapping stated at the head of Appendix A.
\item
  \S{}2.4/\S{}5.8: \texttt{EvenCornerQL} does not use \texttt{hall}; why layer 6 could not see it.
\item
  \S{}3: toolchain \texttt{leanprover/\allowbreak lean4\allowbreak :v4.\allowbreak 33.\allowbreak 1} and the Mathlib pin recorded.
\item
  \S{}4/\S{}5/App. B: receipt and wording corrections --- the batch log's mutation/contradiction
  section cut off at its traceback (the log itself continues to \texttt{LEANCHECK\ DONE}), linter
  cited to \texttt{REPORT\_\allowbreak L10.\allowbreak md} with the 3+3 failure split, 525 decisive kills reported
  separately from 95 timeouts, the bridge file's post-battery date, the scratch-only byte identity,
  \texttt{SCOUT.\allowbreak md} \S{}0 row 18, Clifford's dates aligned to 1870--71/1882, ``forced to be lines'' marked as
  explanation, and the corner's \texttt{V\ =\ 0} derivation cited to \texttt{parent\_\allowbreak crossVec\_\allowbreak zero\_\allowbreak of\_\allowbreak mid}.
\end{enumerate}

\hypertarget{v0.4-independent-review-of-v0.3-review_gpt-6-astra_v03.md-dispositions-in-review_triage_v03b.md}{%
\subsection{\texorpdfstring{v0.4 --- independent review of v0.3 (\texttt{REVIEW\_\allowbreak gpt-6-astra\_\allowbreak v03.\allowbreak md}; dispositions in \texttt{REVIEW\_\allowbreak TRIAGE\_\allowbreak v03b.\allowbreak md})}{v0.4 --- independent review of v0.3 (REVIEW\_gpt-6-astra\_v03.md; dispositions in REVIEW\_TRIAGE\_v03b.md)}}\label{v0.4-independent-review-of-v0.3-review_gpt-6-astra_v03.md-dispositions-in-review_triage_v03b.md}}

\begin{enumerate}
\def\labelenumi{\arabic{enumi}.}
\tightlist
\item
  \S{}2.1: the \texttt{Counter6} ``zero parent cross vector'' example was \textbf{false} and is replaced by the
  \S{}2.4 parabola tuple, with the exact checks named; the \texttt{Counter6} last coordinate is recorded as
  \texttt{−238878720000}.
\item
  \S{}2.1: the piecewise entry formula of \texttt{KP} is displayed; \texttt{Kf}, \texttt{aOf}, \texttt{tOf} are given with lines;
  the kernel equivalence is qualified by nonzero centres, with the \texttt{n\ =\ 1} counterexample noted.
\item
  \S{}2.3: the \texttt{P\_\{k+2\}\ =\ 0} branch now says the \textbf{finite} child points coincide, that zero middle
  coordinates are \texttt{∞}, and that \texttt{heven\_\allowbreak of\_\allowbreak common} supplies the line; no non-uniqueness inference.
\item
  \S{}5.2/\S{}5.4: the two stale references (the ``0/0 row'', the ``definite article in \S{}2.5'') now point at
  the former wording and the battery report's row.
\item
  \S{}5.5: five recorded ssh attempts, not six.
\item
  \S{}1.1: corner-witness margins reported per solver run; the JSON's location in the artifact tree
  stated (the flat referee packet did not carry \texttt{layer0/\allowbreak }).
\item
  \S{}5.3 rewritten for the seven-row layer-13 comparison; \S{}5.11's per-layer provenance updated.
\item
  App. A/B path claims: paper-relative paths and the flat-packet mapping stated.
\item
  \S{}3/App. B: gate count refreshed to 143 occurrences / 98 distinct (three new anchors this round);
  \texttt{L11\_\allowbreak NAMEGATE.\allowbreak md} regenerated.
\end{enumerate}

\hypertarget{v0.5-the-late-block-mutation-run-2026-11-25}{%
\subsection*{v0. --- the late-block mutation run (2026-11-25)}\label{v0.5-the-late-block-mutation-run-2026-11-25}}

No statement, proof or definition changed in this version; the evidence did.

\begin{enumerate}
\def\labelenumi{\arabic{enumi}.}
\tightlist
\item
  \S{}4/\S{}5.1/App. B: mutation coverage of \texttt{Clifford/\allowbreak MorleyKernel.\allowbreak lean} rises from \textbf{61 of 343}
  declarations (all in lines 33--806) to \textbf{318 of 343} (93\%). A late block, run on jack
  2026-11-24 21:28Z $\to$ 2026-11-25 03:41Z, adds 891 mutant rows over the 257 declarations at lines
  816--5586: 872 killed, 3 killed by heartbeat timeout, 16 survivors, 0 unattributed. The round 19--26
  material --- \texttt{heven\_\allowbreak kernel}, \texttt{heven\_\allowbreak line\_\allowbreak of\_\allowbreak P0}, \texttt{evenCorner\_\allowbreak of\_\allowbreak QL}, \texttt{heven\_\allowbreak inf\_\allowbreak all} --- now carries
  mutation rows, so the earlier sentence that no round 19--26 declaration had been mutated is
  withdrawn. Contradiction probes were \textbf{not} extended (layer 5 was skipped), and that stays OPEN
  in v0.5 (closed in v0.6).
\item
  \S{}5.1: the 16 late-block survivors are itemised: nine caught downstream, three vacuous
  (\texttt{2*0+5\ →\ 3*0+5} at \texttt{k\ =\ 0}), and four mutations of \texttt{child\_\allowbreak line\_\allowbreak of\_\allowbreak parent\_\allowbreak mid\_\allowbreak zero}'s \texttt{u\ ≠\ 0}.
  None is a hole.
\item
  \S{}5.8: \texttt{child\_\allowbreak line\_\allowbreak of\_\allowbreak parent\_\allowbreak mid\_\allowbreak zero}'s \texttt{u\ ≠\ 0} added to the unused-hypothesis list.
\item
  \S{}5.5/\S{}4/App. B: john came back online on 2026-11-25 and his like-for-like rerun finished the same
  day --- 1,083 mutants over 324 declarations of \texttt{MorleyKernel.\allowbreak lean}, 1,058 killed, 9 timeout,
  \textbf{16 survivors}, 343 contradiction probes with 0 CONTRADICTORY. The survivor set is the \textbf{same
  sixteen pairs} as jack's, so the two-machine receipt is now claimed for layers 4--5 and \S{}5.5
  rewritten accordingly (\texttt{layer0/\allowbreak john\_\allowbreak receipt\_\allowbreak 20261119/\allowbreak RECEIPT\_\allowbreak VERDICT.\allowbreak md}); the v0.4 snapshot
  figures are marked superseded. The earlier ``unreachable since 2026-11-19'' wording is corrected.
\end{enumerate}

\hypertarget{v0.6-the-late-block-contradiction-probes-and-the-unattributed-rows-2026-11-25}{%
\subsection*{v0. --- the late-block contradiction probes and the unattributed rows (2026-11-25)}\label{v0.6-the-late-block-contradiction-probes-and-the-unattributed-rows-2026-11-25}}

No statement, proof or definition changed in this version either; the evidence did.

\begin{enumerate}
\def\labelenumi{\arabic{enumi}.}
\tightlist
\item
  \S{}4/\S{}5.1/App. B: contradiction probes rise from \textbf{294} declarations to \textbf{551 of the 577}
  (95.4\%). A late block, run on jack 2026-11-25 17:31Z $\to$ 19:09Z over the 257 \texttt{MorleyKernel.\allowbreak lean}
  declarations at lines 816--5592, returns 163 \texttt{ok}, 93 with no hypothesis to probe, 1
  \texttt{probe-invalid} and \textbf{0 CONTRADICTORY}. Layer 5 now covers both blocks of the file that layer 4
  covers, and \S{}5.1 lists the 26 declarations that carry no probe row.
\item
  \S{}5.1: the \texttt{probe-invalid} row is identified --- \texttt{subfamily\_\allowbreak P0\_\allowbreak zero\_\allowbreak gen}, where the \texttt{True} control
  itself fails to compile because \texttt{conj} is bound as an implicit variable under \texttt{autoImplicit} --- and
  is excluded from the \texttt{ok} count rather than reported as a pass.
\item
  \S{}5.1 item 7: the 16 \texttt{error-unattributed} rows were re-run under a lighter load and are all
  decisive, so the battery file has \textbf{935 rows and zero unattributed rows}; \S{}4's layer-4 and
  layer-5 rows and App. B quote the re-run tallies (535 kills, 147 \texttt{ok}) instead of the published
  plus unattributed sums.
\end{enumerate}

\hypertarget{v0.7-the-two-machine-satisfiability-passes-2026-11-26}{%
\subsection*{v0. --- the two-machine satisfiability passes (2026-11-26)}\label{v0.7-the-two-machine-satisfiability-passes-2026-11-26}}

No statement, proof or definition changed in this version; the evidence did.

\begin{enumerate}
\def\labelenumi{\arabic{enumi}.}
\tightlist
\item
  \S{}4/\S{}5.2/App. B: layer 8's ``not performed'' row is replaced by the real record. jack certified 400 of
  596 declarations (67.1\%, 13 h 04 min) and john 394 of 584 (67.5\%, 10 h 17 min) with the same
  instrument at the same pin; \texttt{stmt-invalid} is 0 on both. Across the 577 declarations both machines
  ran, 560 agree, 17 differ only in which machine found a witness, and 397 (68.8\%) are certified by
  at least one machine; 180 are unresolved on both. The section now states precisely what a
  certificate means (an accepted proof that the hypotheses are simultaneously satisfiable), that
  \texttt{unresolved} is a failed search and not a refutation, and that layer 5's \texttt{ok} is neither.
\item
  \S{}5.2: the curated witness table (107 names, jack, 2026-11-26) is reported as having closed none,
  with the reason --- the instrument cannot discharge the propositional side conditions of the
  hypotheses (complex arithmetic; \texttt{decide} cannot close equalities over \texttt{ℂ}) --- and what it would
  need (a per-family proof library for \texttt{Transversal}, \texttt{GenTriple}, \texttt{GenPos4/\allowbreak 5}, \texttt{Chain\ 4\ !{[}\ldots{}{]}},
  \texttt{mD\ ≠\ 0}).
\item
  \S{}4/App. B: layer 8b is recorded as \textbf{0 links} rather than ``no recorded result'': neither tree
  contains a hand-written \texttt{*\_\allowbreak sat} witness.
\end{enumerate}

\hypertarget{v0.8-corrections-from-the-referee-pass-on-v0.7-2026-11-26}{%
\subsection*{v0. --- corrections from the referee pass on v0.7 (2026-11-26)}\label{v0.8-corrections-from-the-referee-pass-on-v0.7-2026-11-26}}

No statement, proof, definition or number changed in this version; nine sentences were corrected
because they said something the record does not support, and three stale provenance statements were
brought up to date. Source: \texttt{REVIEW\_\allowbreak v07\_\allowbreak referee\_\allowbreak bigpickle.\allowbreak md} (one independent instrument, draft-only
input; two further instruments returned no report and are recorded as failed runs) and the
dispositions in \texttt{REVIEW\_\allowbreak TRIAGE\_\allowbreak v07.\allowbreak md}.

\begin{enumerate}
\def\labelenumi{\arabic{enumi}.}
\tightlist
\item
  \S{}1: the novelty sentence now carries its scope in the same breath --- the 26 search rows of
  \texttt{SCOUT.\allowbreak md} \S{}0 and the sources \S{}5.10 records as not searched --- instead of leaving the qualification
  three sections away.
\item
  \S{}1.1: ``The corner is not empty'' became ``not vacuous in practice'', with the numerical search
  downgraded to approximate and exact membership explicitly not claimed from it (the theorem does not
  need it).
\item
  \S{}5.2: the certificate is described correctly --- the hypothesis is a \texttt{Prop}-valued \emph{binder} the
  certificate must inhabit, not an assumption --- and vacuity is no longer inferred from
  \texttt{stmt-invalid\ =\ 0}; 397 witnesses are claimed, 180 declarations are left open.
\item
  \S{}5.3: S1 is removed from the four ``definitional'' rows and named as definitional \textbf{plus} a
  hypothesis-strength gap, which is why it stays open; S4/S5/S6/S7 remain definitional.
\item
  \S{}5.11: ``Nothing in this appendix was run twice on two machines'' was false and is replaced by a
  per-layer statement of which layers ran on both machines; the layer-11 sentence no longer describes
  v0.3 as the current text; the OPEN list no longer names the bridge layer or the pin (both done) and
  names the upload only.
\item
  \S{}3: the tree is no longer described as ``a plain directory, not a repository'' --- it is a git
  repository at \texttt{d15464f} with no remote, and the per-file sha256 table is stated as the identity of
  the audited tree until it is pushed.
\item
  App. A: the \texttt{L11\_\allowbreak NAMEGATE.\allowbreak md} label no longer says v0.3; App. B row 11 records re-runs through
  v0.7; App. C's heading covers v0.3--v0.9; \S{}1.2 separates the polytope citation (Longuet-Higgins

  \begin{enumerate}
  \def\labelenumii{\arabic{enumii})}
  \setcounter{enumii}{1971}
  \tightlist
  \item
    from the cross-ratio route (Bobenko--Schief, arXiv:1509.04109, 2016).
  \end{enumerate}
\end{enumerate}

\hypertarget{v0.9-publication-of-the-code-repository-2026-11-26}{%
\subsection*{v0. --- publication of the code repository (2026-11-26)}\label{v0.9-publication-of-the-code-repository-2026-11-26}}

No mathematical statement, proof, definition or citation changed; this version records the publication
of the artifacts the earlier versions pinned, and corrects the inventory count.

\begin{enumerate}
\def\labelenumi{\arabic{enumi}.}
\tightlist
\item
  \S{}3: the repository is named with its URL and the pin, and the file count is corrected from 153 to
  156 --- commit \texttt{d15464f} also carries \texttt{.\allowbreak gitignore} and the two \texttt{.\allowbreak bak-r26} files, added after the
  earlier count was taken.
\item
  \S{}3: a new \textbf{Availability} paragraph names the repository, the pin, the two check files, the three
  commands that reproduce the gate and the tamper suite, and the \texttt{PaperL} folder of
  \texttt{LEAN-verification}, whose layout is mapped from the artifact-relative paths of Appendix A.
\item
  \S{}5.11: the upload leaves the OPEN list --- the repository is public at the pinned commit and the
  supporting material is published; the Zenodo deposit is recorded there as the authors' remaining
  step, not as done.
\item
  App. B: rows 14 and 15 now record the v0.7 re-cuts (\texttt{L14\_\allowbreak RESTATEMENTS\_\allowbreak v07.\allowbreak md},
  \texttt{L15\_\allowbreak PRIOR\_\allowbreak ART\_\allowbreak v07.\allowbreak md}) that \S{}5.11 and \S{}5 already described, and row 11's re-run list is brought
  up to v0.9 with a typo fixed.
\end{enumerate}

\hypertarget{v0.10-corrections-from-the-referee-pass-on-v0.9-2026-11-26}{%
\subsection*{v0. --- corrections from the referee pass on v0.9 (2026-11-26)}\label{v0.10-corrections-from-the-referee-pass-on-v0.9-2026-11-26}}

Source: \texttt{REVIEW\_\allowbreak v09\_\allowbreak astra.\allowbreak md} (one independent instrument; input: this draft, the pinned tree and the
published reports), dispositions in \texttt{REVIEW\_\allowbreak TRIAGE\_\allowbreak v09.\allowbreak md}. Two findings were MAJOR and eight MINOR;
all ten were accepted after each was checked against the source or the record. No Lean proof,
statement or definition changed.

\begin{enumerate}
\def\labelenumi{\arabic{enumi}.}
\tightlist
\item
  \S{}5.3: the layer-13 classification of S4/S5 is corrected. They are not definitional variants: the
  writer's finite-plane encoding forbids $\infty$ at every level, so a level object that is a line is
  unrepresentable. The recorded \texttt{BlindGenPos} triple with collinear second points satisfies the
  writer's hypotheses and admits no blind finite-plane level-3 object, so the writer's S4 statement
  is false as written at \texttt{n\ =\ 3}; the library's cutoff includes \texttt{n\ =\ 3} and is unaffected. S6/S7
  remain definitional, and both directions are bridged.
\item
  \S{}5.1/\S{}5.8/App. C: the four surviving mutants of \texttt{child\_\allowbreak line\_\allowbreak of\_\allowbreak parent\_\allowbreak mid\_\allowbreak zero} are re-described ---
  they alter the \textbf{conclusion} of a lemma whose hypotheses are inconsistent on purpose, so the
  branch is empty and no proof obligation reaches them. (v0.11 refines this further: only two of the
  four touch the \texttt{u\ ≠\ 0} conjunct; one flips the finiteness antecedent and one the incidence
  equation.) The lemma leaves the unused-input account; the layer-6 501/8 tally is unchanged.
\item
  \S{}3: the tamper suite is \textbf{eight} tests, not nine, and the reproducer is \texttt{bash\ tests/\allowbreak tamper.\allowbreak sh}
  (the file is mode 100644 at the pin); the citation gate makes \textbf{25} name-attribution checks, not

  \begin{enumerate}
  \def\labelenumii{\arabic{enumii}.}
  \setcounter{enumii}{21}
  \tightlist
  \item
  \end{enumerate}
\item
  \S{}3/App. A: the supporting folder is described by its layout instead of as a mirror, the
  artifact-relative path convention is stated once, and john's original battery report is removed
  from the externally-hosted class --- it is published in \texttt{layer0/\allowbreak john\_\allowbreak receipt\_\allowbreak 20261119/\allowbreak }.
\item
  \S{}1.1: the rank-route sentence now says numerical evidence, not refutation, and ties the \texttt{1e-9}
  residual and the \texttt{1e-8} collinearity figure to the earlier, looser run (\texttt{3.\allowbreak 9e-5} collinearity for
  the stored tuple).
\item
  \S{}5.1/\S{}5.5/App. C: the stale satisfiability OPEN item is removed (only the expected-job manifest
  remains), the round-26b timings are labelled by machine and operation, and the v0.9 entry no longer
  says that no number changed.
\item
  \S{}4: layer 5's warning now says that a failed search does not certify \textbf{satisfiability}.
\item
  \S{}5.1/\S{}5.11/App. B: ``checks alone'' and the machine-count summaries are scoped --- layers 4/5 on john
  cover \texttt{MorleyKernel.\allowbreak lean}, layer 13 on john covers \texttt{BlindCompareFull.\allowbreak lean}, and the bridge file is
  jack-only.
\item
  Supporting reports: \texttt{L13\_\allowbreak BRIDGE\_\allowbreak STATUS.\allowbreak md}'s stale \texttt{sorryAx} line and \texttt{L13\_\allowbreak FULL\_\allowbreak COMPARE.\allowbreak md}'s
  ``only definitional row'' sentence are corrected with a dated note, and \texttt{L13\_\allowbreak REPORT.\allowbreak md} carries a
  supersession header. The audited pin and the raw logs are untouched.
\end{enumerate}

\hypertarget{v0.11-corrections-from-the-opus-5.5-pass-on-v0.10-2026-11-26}{%
\subsection*{v0. --- corrections from the Opus 5.5 pass on v0.10 (2026-11-26)}\label{v0.11-corrections-from-the-opus-5.5-pass-on-v0.10-2026-11-26}}

Source: \texttt{REVIEW\_\allowbreak v10\_\allowbreak opus.\allowbreak md} (one independent instrument; inputs: this draft, the pinned tree, the
published repositories; no compilation), dispositions in \texttt{REVIEW\_\allowbreak TRIAGE\_\allowbreak v10.\allowbreak md}. One finding was
MAJOR, ten MINOR and two NOTE; all thirteen were accepted after each was checked against the data or
the source. No Lean proof, statement or definition changed.

\begin{enumerate}
\def\labelenumi{\arabic{enumi}.}
\tightlist
\item
  \S{}5.2: layer 8's 397 certificates are now split by statement class --- 210 hypothesis-free, 119
  closed, 68 inhabiting a hypothesis --- so the paper no longer says that 397 ``carry a witness of
  their hypotheses''. It also records that a hypothesis packed inside a \texttt{Prop} definition is
  invisible to layer 8 as it is to layer 6, and that \texttt{parent\_\allowbreak crossVec\_\allowbreak zero\_\allowbreak of\_\allowbreak mid} is unresolved on
  both machines, so whether the corner is inhabited rests on \S{}1.1's numerical witness.
\item
  \S{}5.1/\S{}5.8: only two of the four \texttt{child\_\allowbreak line\_\allowbreak of\_\allowbreak parent\_\allowbreak mid\_\allowbreak zero} survivors touch the \texttt{u\ ≠\ 0}
  conjunct; one flips the grandchild-finiteness antecedent and one the incidence equation.
\item
  \S{}5.3: S6 alone is definitional; S7 is definitional \textbf{plus} a hypothesis-strength gap, closed
  because the writer's statement is proved outright. S4's falsity is labelled a hand argument on
  exact coordinates, not a machine-checked statement. The S1 paragraph records the direct route in
  \texttt{layer13/\allowbreak BlindChain.\allowbreak lean} (levels 3--4 for every \texttt{BlindGenPos} tuple, level 5 conditional, all-\texttt{n}
  reduced to \texttt{blindClifford\_\allowbreak second}) and no longer offers rescoping, which the recorded note rules
  out.
\item
  \S{}5.1: the 19 un-mutated declarations are ones on which no mutant was generated, not ones the
  generator ``cannot'' alter; the headline is no longer called a corollary; and john's receipt is
  recorded as probing all 343 \texttt{MorleyKernel} declarations, the 26 jack left unprobed included.
\item
  \S{}4/\S{}5.1/App. B/App. C: the late-block probe range is 816--5586 --- \texttt{heven\_\allowbreak inf\_\allowbreak all} is the last probed
  declaration --- not 5592.
\item
  \S{}3: the citation gate leaves eight ``lines A--B'' range tokens (four distinct spans), and the
  reproduction steps add \texttt{lake\ exe\ cache\ get} and say the manifest, not the lakefile tag, fixes
  Mathlib.
\item
  App. A/\S{}5.11/App. B: john's battery report is published and layers 1, 2 and 7 are not jack-only;
  the three layer-11 provenance statements now agree (through v0.11); the v0.9$\to$v0.10 and
  v0.10$\to$v0.11 review rounds are recorded; and App. A gains a table mapping every artifact path to
  its published location.
\item
  \S{}1.1: the numerical-witness caveat is stated once and the duplicated sentence is merged.
\item
  Supporting reports: the three reports that still described \texttt{u\ ≠\ 0} as an unused hypothesis carry a
  dated correction notice.
\end{enumerate}

\hypertarget{v0.12-typeset-in-the-style-of-papers-i-and-k-2026-11-26}{%
\subsection*{v0. --- typeset in the style of Papers I and K (2026-11-26)}\label{v0.12-typeset-in-the-style-of-papers-i-and-k-2026-11-26}}

No statement, proof, number or citation changed. The paper is now built as LaTeX with the preamble
used by the other papers in this series (11 pt article, 1 in margins, Latin Modern, booktabs
tables, numbered sections, title/author block with the licence footnote), instead of the markdown
pandoc build. The abstract, previously only in the deposit record, is now part of the paper. The
version changelog (this appendix) stays in the markdown source and is not typeset into the PDF;
\texttt{paperl.\allowbreak tex} is generated from \texttt{DRAFT.\allowbreak md} by \texttt{build\_\allowbreak paperl.\allowbreak py}, and both are published beside the
PDF.

\end{document}
